\documentclass[draft,linenumbers]{agujournal2019}
\usepackage{url}
\usepackage{lineno}
\usepackage{soul}

\usepackage{xcolor}
\definecolor{default-linkcolor}{HTML}{A50000}
\definecolor{default-filecolor}{HTML}{A50000}
\definecolor{default-citecolor}{HTML}{4077C0}
\definecolor{default-urlcolor}{HTML}{4077C0}

\usepackage{hyperref}
\hypersetup{
  hidelinks,
  breaklinks=true
}


\providecommand{\tightlist}{\setlength{\itemsep}{0pt}\setlength{\parskip}{0pt}}
\providecommand{\listfigurename}{List of Figures}
\providecommand{\listtablename}{List of Tables}
\newcounter{none}
\setcounter{none}{0}

\makeatletter
\@ifundefined{pandocbounded}{
  \usepackage{graphicx}
  \newcommand{\pandocbounded}[1]{\setbox0=\hbox{#1}\ifdim\wd0>\linewidth\resizebox{\linewidth}{!}{#1}\else#1\fi}
}{}
\makeatother

\usepackage{longtable,booktabs,array}
\usepackage{calc}
\usepackage{etoolbox}
\makeatletter
\patchcmd\longtable{\par}{\if@noskipsec\mbox{}\fi\par}{}{}
\makeatother
\IfFileExists{footnotehyper.sty}{\usepackage{footnotehyper}}{\usepackage{footnote}}
\makesavenoteenv{longtable}

\NewDocumentCommand\citeproctext{}{}
\NewDocumentCommand\citeproc{mm}{\begingroup\def\citeproctext{#2}\cite{#1}\endgroup}
\makeatletter
 \let\@cite@ofmt\@firstofone
 \def\@biblabel#1{}
 \def\@cite#1#2{{#1\if@tempswa , #2\fi}}
\makeatother

\newlength{\cslhangindent}
\setlength{\cslhangindent}{1.5em}
\newlength{\csllabelwidth}
\setlength{\csllabelwidth}{3em}
\newenvironment{CSLReferences}[2]
 {\begin{list}{}{
  \setlength{\itemindent}{0pt}
  \setlength{\leftmargin}{0pt}
  \setlength{\parsep}{0pt}
  \ifodd #1
   \setlength{\leftmargin}{\cslhangindent}
   \setlength{\itemindent}{-\cslhangindent}
  \fi
  \setlength{\itemsep}{#2\baselineskip}
  }
  \item\relax}
 {\end{list}}
\usepackage{calc}
\newcommand{\CSLBlock}[1]{#1\hfill\break}
\newcommand{\CSLLeftMargin}[1]{\parbox[t]{\csllabelwidth}{#1}}
\newcommand{\CSLRightInline}[1]{\parbox[t]{\linewidth - \csllabelwidth}{#1}\break}
\newcommand{\CSLIndent}[1]{\hspace{\cslhangindent}#1}

\usepackage{amsmath}
\usepackage{amssymb}
\usepackage{amsfonts}
\makeatletter
\@ifpackageloaded{subcaption}{}{\usepackage{subcaption}}
\@ifpackageloaded{caption}{}{\usepackage{caption}}
\captionsetup[subfigure]{margin=0.5em}
\AtBeginDocument{%
\renewcommand*\figurename{Figure}
\renewcommand*\tablename{Table}
}
\AtBeginDocument{%
\renewcommand*\listfigurename{List of Figures}
\renewcommand*\listtablename{List of Tables}
}
\newcounter{pandoccrossref@subfigures@footnote@counter}
\newenvironment{pandoccrossrefsubfigures}{%
\setcounter{pandoccrossref@subfigures@footnote@counter}{0}
\begin{figure}\centering%
\gdef\global@pandoccrossref@subfigures@footnotes{}%
\DeclareRobustCommand{\footnote}[1]{\footnotemark%
\stepcounter{pandoccrossref@subfigures@footnote@counter}%
\ifx\global@pandoccrossref@subfigures@footnotes\empty%
\gdef\global@pandoccrossref@subfigures@footnotes{{##1}}%
\else%
\g@addto@macro\global@pandoccrossref@subfigures@footnotes{, {##1}}%
\fi}}%
{\end{figure}%
\addtocounter{footnote}{-\value{pandoccrossref@subfigures@footnote@counter}}
\@for\f:=\global@pandoccrossref@subfigures@footnotes\do{\stepcounter{footnote}\footnotetext{\f}}%
\gdef\global@pandoccrossref@subfigures@footnotes{}}
\@ifpackageloaded{float}{}{\usepackage{float}}
\floatstyle{ruled}
\@ifundefined{c@chapter}{\newfloat{codelisting}{h}{lop}}{\newfloat{codelisting}{h}{lop}[chapter]}
\floatname{codelisting}{Listing}
\newcommand*\listoflistings{\listof{codelisting}{List of Listings}}
\makeatother

\journalname{JGR: Solid Earth}

\begin{document}


\title{Beyond Equilibrium II: Ringwoodite Decomposition Kinetics Amplify
660 km Discontinuity Depressions and Offer an Independent Control on
Slab Stagnation}

\authors{Buchanan Kerswell\affil{1}, John Wheeler\affil{1}, Rene
Gassmöller\affil{2}, J. Huw Davies\affil{3}, Isabel
Papanagnou\affil{4}, Sanne Cottaar\affil{4}, David P. Dobson\affil{5}}

\affiliation{1}{Department of Earth, Oceans and Ecological Sciences,
University of Liverpool, Liverpool L69 3BX, UK}
\affiliation{2}{GEOMAR Helmholtz Centre for Ocean Research Kiel, Kiel,
Germany}
\affiliation{3}{School of Earth and Environmental Sciences, Cardiff
University, Park Place, Cardiff, CF10 3AT, UK}
\affiliation{4}{Bullard Laboratories, Department of Earth Sciences,
University of Cambridge, Cambridge CB3 0EZ, UK}
\affiliation{5}{Department of Earth Sciences, University College London,
London WC1E 6BS, UK}

\correspondingauthor{Buchanan Kerswell}{b.kerswell@liverpool.ac.uk}

\begin{keypoints}
\item Plumes consistently uplift the 660; slabs show kinetically
controlled deepening and broadening across three distinct regimes
\item Diffusion-controlled kinetics in cold slabs produce monotonically
broader and deeper 660s as reaction rates slow
\item Ringwoodite kinetics amplify 660 depressions and reduce slabs
velocities, partially explaining slab stagnation and 660 topography
\end{keypoints}

\begin{abstract}
The 660 km seismic discontinuity (`660') varies beneath subduction
zones, ranging from modest depressions consistent with equilibrium
thermodynamics to anomalously deep, broad transitions challenging purely
thermal interpretations. Laboratory experiments establish that the
ringwoodite decomposition responsible for the 660 proceeds via
diffusion-controlled kinetics, yet their quantitative effects on 660
structure and slab dynamics have not been evaluated in compressible
mantle flow simulations. We couple ringwoodite decomposition kinetics to
simulations of mantle plumes and subducting slabs, varying kinetic
parameters across six orders of magnitude. Results reveal a fundamental
asymmetry between hot and cold environments. In plumes, elevated
temperatures produce uplifted, narrow 660s that are relatively
insensitive to kinetics due to rapid reaction rates. In slabs, kinetics
exert strong control through three regimes: 1) a quasi-equilibrium
regime producing 660s consistent with Clapeyron-slope predictions; 2) an
intermediate regime where metastable ringwoodite accumulates within a
broadening buoyant zone, slowing slab descent; and 3) a stagnation
regime where slabs pond above the 660, producing the deepest, broadest
660s and slowest descent velocities. The endothermic Clapeyron slope and
kinetic inhibition reinforce each other, amplifying temperature-driven
660 deflection by a factor of 2--4 while independently promoting slab
stagnation. Slabs smoothly transition through each regime as reaction
rates slow, without abrupt reactions at overstepped metastable
conditions. These findings establish ringwoodite kinetics as a
first-order control on slab stagnation dynamics and 660 seismic
structure under Earth-like conditions, explaining the diverse slab
behaviors observed at the base of the mantle transition zone.
\end{abstract}

\section*{Plain Language Summary}
Earth has a boundary 660 kilometers underground where the mineral
ringwoodite transforms into two new minerals. Beneath sinking tectonic
plates, this boundary is often deeper and thicker than expected.
Scientists assumed cold sinking plates were the main cause, but
laboratory experiments show the ringwoodite transformation is also
limited by how quickly silicon atoms move through rock---a speed limit
whose effect on mantle flow was untested. We used computer simulations
of rising mantle plumes and sinking slabs to test how reaction speed
affects slab motion and the seismic appearance of the 660 km boundary.
In hot plumes, high temperatures allow rapid transformation, so the 660
shifts slightly upward and remains narrow regardless of reaction speed.
In cold slabs, however, reaction speed strongly affects flow behavior.
Fast reactions allow slabs to sink steadily through the 660 with modest
deepening. Slower reactions cause untransformed ringwoodite to build up
inside the slab, slowing its descent while pushing the 660 deeper and
wider. When reactions are very slow, slabs can stall above the 660
instead of sinking into the lower mantle, producing the deepest,
broadest seismic signals without abrupt re-sharpening. Our results show
that ringwoodite transformation speed is a major control on tectonic
plate movement through Earth's interior. It can amplify
temperature-driven boundary deflection by a factor of 2 to 4 and slow
slab descent by up to ten times, explaining why some slabs sink into the
lower mantle while others flatten and temporarily stagnate above the
transition zone.

\clearpage

\section*{Definition of Symbols}\label{sec:definition-of-symbols}
\addcontentsline{toc}{section}{Definition of Symbols}

{\def\LTcaptype{none} % do not increment counter
\begin{longtable}[]{@{}
  >{\raggedright\arraybackslash}p{(\linewidth - 6\tabcolsep) * \real{0.4688}}
  >{\raggedright\arraybackslash}p{(\linewidth - 6\tabcolsep) * \real{0.2188}}
  >{\raggedright\arraybackslash}p{(\linewidth - 6\tabcolsep) * \real{0.1562}}
  >{\raggedright\arraybackslash}p{(\linewidth - 6\tabcolsep) * \real{0.1562}}@{}}
\toprule\noalign{}
\begin{minipage}[b]{\linewidth}\raggedright
Parameter
\end{minipage} & \begin{minipage}[b]{\linewidth}\raggedright
Symbol
\end{minipage} & \begin{minipage}[b]{\linewidth}\raggedright
Unit
\end{minipage} & \begin{minipage}[b]{\linewidth}\raggedright
Equations
\end{minipage} \\
\midrule\noalign{}
\endhead
\bottomrule\noalign{}
\endlastfoot
Activation energy & \(E^{\ast}\) & J mol\(^{-1}\) &
\ref{eq:growth-rate-660}, \ref{eq:reaction-rate-660} \\
Activation enthalpy & \(H^{\ast}\) & J mol\(^{-1}\) &
\ref{eq:reaction-rate-410} \\
Activation volume & \(V^{\ast}\) & m\(^3\) mol\(^{-1}\) &
\ref{eq:reaction-rate-410} \\
Activation factor (rheology) & \(B\) & - & \ref{eq:rheological-model} \\
Compressibility (reference) & \(\bar{\beta}\) & Pa\(^{-1}\) &
\ref{eq:density-ala} \\
Density & \(\rho\) & kg m\(^{-3}\) &
\ref{eq:navier-stokes-no-inertia}, \ref{eq:continuity-compressible}, \ref{eq:energy}, \ref{eq:continuity-expanded}, \ref{eq:density-ala} \\
Density (reference) & \(\bar{\rho}\) & kg m\(^{-3}\) &
\ref{eq:adiabatic-pressure}, \ref{eq:density-ala} \\
Density (dynamic) & \(\hat{\rho}\) & kg m\(^{-3}\) & - \\
Deviatoric stress tensor & \(\sigma^{\prime}\) & Pa &
\ref{eq:navier-stokes-no-inertia}, \ref{eq:energy} \\
Deviatoric strain rate tensor & \(\dot{\epsilon}^{\prime}\) & s\(^{-1}\)
& \ref{eq:energy} \\
Gas constant & \(R\) & J mol\(^{-1}\) K\(^{-1}\) &
\ref{eq:growth-rate-660}, \ref{eq:rheological-model} \\
Gibbs width factor & \(w\) & J mol\(^{-1}\) &
\ref{eq:transition-function} \\
Grain size & \(d\) & m & - \\
Gravitational acceleration & \(g\) & m s\(^{-2}\) &
\ref{eq:navier-stokes-no-inertia}, \ref{eq:adiabatic-temperature}, \ref{eq:adiabatic-pressure} \\
Growth rate & \(\dot{x}\) & m s\(^{-1}\) &
\ref{eq:volume-fraction}, \ref{eq:volume-transformation-rate}, \ref{eq:growth-rate-660} \\
Growth rate prefactor & \(K^{\prime}\) & m\(^7\) J\(^{-2}\) s\(^{-1}\) &
\ref{eq:growth-rate-660} \\
Kinetic prefactor (olivine \(\Leftrightarrow\) wadsleyite) &
\(Z_\mathrm{ol}\) & K s\(^{-1}\) & \ref{eq:reaction-rate-410} \\
Kinetic prefactor (wadsleyite \(\Leftrightarrow\) ringwoodite) &
\(Z_\mathrm{wd}\) & K s\(^{-1}\) & \ref{eq:reaction-rate-410} \\
Kinetic prefactor (ringwoodite \(\Leftrightarrow\) bridgmanite +
periclase) & \(Z_\mathrm{ri}\) & mol\(^2\) J\(^{-2}\) s\(^{-1}\) &
\ref{eq:reaction-rate-660} \\
Latent heat & \(Q_L\) & J kg\(^{-1}\) & \ref{eq:energy} \\
Molar entropy & \(\bar{S}\) & J mol\(^{-1}\) K\(^{-1}\) &
\ref{eq:gibbs} \\
Molar Gibbs free energy & \(\bar{G}\) & J mol\(^{-1}\) &
\ref{eq:gibbs}, \ref{eq:transition-function} \\
Molar volume & \(\bar{V}\) & m\(^{3}\) mol\(^{-1}\) & \ref{eq:gibbs} \\
Nucleation site factor & \(N\) & m\(^{-1}\) &
\ref{eq:volume-fraction}, \ref{eq:volume-transformation-rate} \\
Pressure & \(P\) & Pa &
\ref{eq:navier-stokes-no-inertia}, \ref{eq:energy} \\
Pressure (reference) & \(\bar{P}\) & Pa & \ref{eq:adiabatic-pressure} \\
Pressure (dynamic) & \(\hat{P}\) & Pa &
\ref{eq:density-ala}, \ref{eq:gibbs} \\
Reaction rate & \(\frac{dX}{dt}\); \(\dot{X}\) & s\(^{-1}\) &
\ref{eq:volume-transformation-rate}, \ref{eq:reaction-rate-410}, \ref{eq:reaction-rate-660}, \ref{eq:composition} \\
Reaction completion timescale & \(\tau\) & Ma & - \\
Reaction completion length scale & \(L\) & km & - \\
Specific heat capacity (reference) & \(\bar{C}_p\) & J kg\(^{-1}\)
K\(^{-1}\) & \ref{eq:energy}, \ref{eq:adiabatic-temperature} \\
Temperature & \(T\) & K &
\ref{eq:energy}, \ref{eq:growth-rate-660}, \ref{eq:rheological-model} \\
Temperature (reference) & \(\bar{T}\) & K &
\ref{eq:adiabatic-temperature}, \ref{eq:rheological-model} \\
Temperature (dynamic) & \(\hat{T}\) & K &
\ref{eq:density-ala}, \ref{eq:gibbs}, \ref{eq:rheological-model} \\
Thermal conductivity (reference) & \(\bar{k}\) & W m\(^{-1}\) K\(^{-1}\)
& \ref{eq:energy} \\
Thermal expansivity (reference) & \(\bar{\alpha}\) & K\(^{-1}\) &
\ref{eq:energy}, \ref{eq:adiabatic-temperature}, \ref{eq:density-ala} \\
Time & \(t\) & s &
\ref{eq:continuity-compressible}, \ref{eq:energy}, \ref{eq:continuity-expanded}, \ref{eq:volume-fraction}, \ref{eq:composition} \\
Transition function & \(f_j\) & - &
\ref{eq:transition-function}, \ref{eq:equilibrium-fractions} \\
Velocity (field) & \(\vec{u}\) & m s\(^{-1}\) &
\ref{eq:continuity-compressible}, \ref{eq:energy}, \ref{eq:continuity-expanded}, \ref{eq:composition} \\
Velocity (boundary) & \(\vec{u}_\mathrm{in}\) & cm yr\(^{-1}\) & - \\
Velocity (maximum through 660) & \(\vec{v}_\mathrm{max}\) & cm
yr\(^{-1}\) & - \\
Viscosity & \(\eta\) & Pa s & \ref{eq:rheological-model} \\
Viscosity (reference) & \(\bar{\eta}\) & Pa s &
\ref{eq:rheological-model} \\
Viscosity jump factor (lower mantle) & \(b\) & - &
\ref{eq:rheological-model}, \ref{eq:rheological-jump-factor} \\
Viscosity prefactor (phase \(i\)) & \(\nu\) & - &
\ref{eq:rheological-jump-factor} \\
Volume fraction (actual) & \(X\) & - &
\ref{eq:volume-fraction}, \ref{eq:volume-transformation-rate}, \ref{eq:reaction-rate-410}, \ref{eq:reaction-rate-660}, \ref{eq:composition} \\
Volume fraction (equilibrium) & \(\phi\) & - &
\ref{eq:equilibrium-fractions}, \ref{eq:rheological-jump-factor} \\
\end{longtable}
}

\clearpage

\section{Introduction}\label{sec:introduction}

Earth's 660 km seismic discontinuity (hereafter the `660') marks the
base of the mantle transition zone (MTZ), where ringwoodite decomposes
into bridgmanite and periclase (the `post-spinel' reaction). Two
shallower discontinuities, near 410 and 520 km depth (hereafter the
`410' and `520'), mark the olivine \(\Leftrightarrow\) wadsleyite and
wadsleyite \(\Leftrightarrow\) ringwoodite transitions that further
define the MTZ. High-pressure experiments constrain the Clapeyron slope
of the post-spinel reaction to -1 to -3 MPa/K ({Ishii et al.}, 2019;
Litasov et al., 2005), which deepens the phase transition beneath cold
subducting slabs and has long been recognized as a potential barrier to
whole-mantle convection (Ita \& Stixrude, 1992; Ringwood, 1991).

Equilibrium thermodynamics explains the 660's first-order behavior, but
two observations suggest it is incomplete. First, 660 structure itself
is more variable than purely thermal effects predict. While global
seismic studies show patterns broadly consistent with the thermal
modulation expected from a negative Clapeyron slope (e.g., Flanagan \&
Shearer, 1998; Lebedev et al., 2002), amplitude variations point to
additional compositional heterogeneity ({Goes et al.}, 2022; Waszek et
al., 2021; {Yu et al.}, 2023), and local 660 structures in many
subduction zones show stronger than predicted deflections (e.g., Cottaar
\& Deuss, 2016; Ba et al., 2025; Tang et al., 2014). Proposed
explanations include water in ringwoodite (Muir et al., 2021), the
akimotoite \(\Leftrightarrow\) bridgmanite transition ({Chanyshev et
al.}, 2022; Cottaar \& Deuss, 2016), and grain-size-dependent viscosity
reduction (Kubo et al., 2000; Mao \& Zhong, 2021). Second, slab dynamics
at the 660 depend on more than the Clapeyron slope alone. Numerical
geodynamic studies show that slab stagnation above the 660 requires a
negative Clapeyron slope, a lower-mantle viscosity increase, and trench
retreat acting jointly (Garel et al., 2014; {Goes et al.}, 2017).
Whether reaction kinetics can independently explain these observations
has not been evaluated. Few studies have moved beyond the equilibrium
paradigm to test how much microscale reaction kinetics influence
macroscale flow dynamics and seismic structures.

Laboratory experiments establish that the 410 olivine
\(\Leftrightarrow\) wadsleyite reaction and the 660 post-spinel reaction
are governed by fundamentally different physical mechanisms (Dobson \&
Mariani, 2014; Hosoya et al., 2005; Kubo et al., 2002; Lessing et al.,
2022), which translate into different semi-empirical kinetic rate
equations with distinct mathematical formulations and parameter
dependencies. In a companion study, Kerswell et al. (2026) coupled
\emph{interface-controlled} olivine \(\Leftrightarrow\) wadsleyite
kinetics to compressible mantle flow simulations and identified three
distinct kinetic regimes in cold slabs with diagnostic seismic
signatures at the 410. Because \emph{diffusion-controlled} ringwoodite
decomposition kinetics takes a different mathematical form, it is
expected to produce qualitatively different dynamic behavior at the 660.
Whether these two reaction mechanisms produce different flow dynamics
and discontinuity structures, and whether reaction kinetics at the 660
matter at all for large-scale mantle flow, has not been systematically
evaluated in geodynamic simulations of Earth's upper mantle.

Here, we couple experimentally calibrated ringwoodite decomposition
kinetics (after Kubo et al., 2002) to compressible mantle flow
simulations using ASPECT (Heister et al., 2017; Kronbichler et al.,
2012). Our simulations simultaneously resolve the 410, 520, and 660
transitions, enabling the first direct comparison between reaction
mechanisms in a consistent dynamic framework. We hypothesize that
ringwoodite decomposition kinetics control whether slabs penetrate or
stagnate at the 660, such that slower reaction rates sustain a buoyant
metastable zone that impedes descent while faster rates facilitate
continuous penetration. We further hypothesize that kinetic inhibition
thickens and flattens slabs above the 660, producing seismically broad,
deeply displaced discontinuities in stagnant settings and narrow,
modestly displaced signals where kinetics are fast. These hypotheses
motivate three specific questions:

\begin{enumerate}
\def\labelenumi{\arabic{enumi}.}
\tightlist
\item
  Do ringwoodite decomposition kinetics control 660 structure and slab
  dynamics, and how do these effects differ from those at the 410?
\item
  How do rheological strength contrasts and slab kinematics modulate
  kinetic effects on 660 structure?
\item
  Can seismic observations of 660 structure independently constrain
  ringwoodite decomposition kinetics?
\end{enumerate}

We vary ringwoodite decomposition kinetics across six orders of
magnitude, spanning the range permitted by typical grain sizes under
nominally dry conditions (Karato, 1984, 2008; Kubo et al., 2002). We
also vary rheological strength contrasts, inflow velocities, and
lower-mantle viscosity contrasts. The results show three kinetic regimes
at the 660 that amplify temperature-driven 660 displacements by a factor
of 2--4 and reduce slab descent rates by up to an order of magnitude.
Unlike the kinetic regimes at the 410, which include an abrupt reaction
threshold at highly-overstepped metastable conditions, the 660 regimes
change monotonically from quasi-equilibrium to stagnation. We discuss
how the differences between the two reaction mechanisms produce
fundamentally different slab dynamics and discontinuity structures, and
how jointly observing the 410 and 660 in the same slab column offers a
seismological pathway toward discriminating kinetic from thermal and
compositional contributions to MTZ structure.

\section{Methods}\label{sec:methods}

\subsection{Governing Equations for Compressible Mantle
Flow}\label{sec:governing-equations-for-compressible-mantle-flow}

We solve the same governing equations for compressible mantle flow as
Kerswell et al. (2026), summarized here for completeness. We simulate
mantle flow using the finite-element geodynamic code ASPECT v3.0.0
(Bangerth et al., 2018, 2024; Clevenger \& Heister, 2021; {Fraters et
al.}, 2019; Fraters, 2020; Gassmöller et al., 2018) to determine the
velocity \(\vec{u}\), pressure \(P\), and temperature \(T\) fields that
satisfy the following equations:

\begin{equation}\protect\phantomsection\label{eq:navier-stokes-no-inertia}{
  \nabla P - \nabla \cdot \sigma^{\prime} = \rho\, g
}\end{equation}

\begin{equation}\protect\phantomsection\label{eq:continuity-compressible}{
  \frac{\partial \rho}{\partial t} + \nabla \cdot (\rho\, \vec{u}) = 0
}\end{equation}

\begin{equation}\protect\phantomsection\label{eq:energy}{
  \rho\, \bar{C}_p \left(\frac{\partial T}{\partial t} + \vec{u} \cdot \nabla T \right) - \nabla \cdot \left(\bar{k}\, \nabla T \right) = \sigma^{\prime} : \dot{\epsilon}^{\prime} + \bar{\alpha}\, T \left(\vec{u} \cdot \nabla P \right) + Q_L
}\end{equation}

where \(\sigma^{\prime}\) is the deviatoric stress tensor,
\(\dot{\epsilon}^{\prime}\) is the deviatoric strain rate tensor,
\(\rho\) is density, \(g\) is gravitational acceleration, \(t\) is time,
and \(\bar{C}_p\), \(\bar{k}\), \(\bar{\alpha}\) are reference specific
heat capacity, thermal conductivity, and thermal expansivity
respectively. The term \(Q_L\) captures latent heat from phase
transitions. We adopt the projected density approximation (Gassmöller et
al., 2020), reformulating Equation~\ref{eq:continuity-compressible} as:

\begin{equation}\protect\phantomsection\label{eq:continuity-expanded}{
  \frac{1}{\rho} \frac{\partial \rho}{\partial t} + \nabla \cdot \vec{u} + \left(\frac{1}{\rho} \nabla \rho \right) \cdot \vec{u} = 0
}\end{equation}

The projected density \(\rho(T, P, X)\) varies with temperature,
pressure, and volume fractions of the relevant mineral phases, ensuring
that local density changes from both pressure-temperature (PT)
variations and phase transitions influence the flow through buoyancy,
volumetric expansion, and compression. See Gassmöller et al. (2020) for
a detailed discussion of this compressible formulation.

\subsection{Numerical Setup}\label{sec:numerical-setup}

\subsubsection{Adiabatic Reference
Conditions}\label{sec:adiabatic-reference-conditions}

We compute adiabatic reference conditions using the same approach as
Kerswell et al. (2026), extended to encompass the full MTZ and uppermost
lower mantle (Figure~\ref{fig:PYR-material-table}). We evaluate entropy
changes over a PT range of 1053--2673 K and 0.001--43 GPa using Gibbs
free energy minimization with Perple\_X (v7.0.9, Connolly, 2009),
assuming a dry pyrolitic bulk composition (Green et al., 1979) in the
Na\(_2\)O-CaO-FeO-MgO-Al\(_2\)O\(_3\)-SiO\(_2\) (NCFMAS) chemical system
with the equations of state and solution models of Stixrude \&
Lithgow-Bertelloni (2022). We determine an isentropic adiabat using the
Newton-Raphson method and evaluate reference material properties
(\(\bar{\rho}\), \(\bar{\alpha}\), \(\bar{C}_p\), \(\bar{\beta}\)) along
the adiabat using BurnMan (Cottaar et al., 2014; Myhill et al., 2023)
with equations of state from Stixrude \& Lithgow-Bertelloni (2022) for
pure Mg ringwoodite, bridgmanite, and periclase (Figure S1).

\begin{figure}
\centering
\includegraphics[width=0.65\linewidth,height=\textheight,keepaspectratio,alt={Entropy (left) and density (right) changes across Earth's MTZ under thermodynamic equilibrium and hydrostatic stress conditions. The black box indicates the approximate PT range within our ASPECT simulations. The white line indicates the isentropic adiabat used to calculate reference material properties.}]{../figs/adiabat/PYR-material-table.png}
\caption{Entropy (left) and density (right) changes across Earth's MTZ
under thermodynamic equilibrium and hydrostatic stress conditions. The
black box indicates the approximate PT range within our ASPECT
simulations. The white line indicates the isentropic adiabat used to
calculate reference material properties.}\label{fig:PYR-material-table}
\end{figure}

\subsubsection{Initialization and Boundary
Conditions}\label{sec:initialization-and-boundary-conditions}

We conduct simulations in a 1500 \(\times\) 1000 km rectangular model
domain initialized with equilibrium mineral assemblages across the
olivine, wadsleyite, ringwoodite, and post-spinel stability fields. An
applied surface pressure of 10 GPa and surface temperature of 1706 K
places the 410, 520, and 660 transitions at approximately the center of
the model domain (Figure~\ref{fig:initial-setup-diagram}). We compute
initial adiabatic PT profiles by numerically integrating:

\begin{equation}\protect\phantomsection\label{eq:adiabatic-temperature}{
  \frac{d\bar{T}}{dy} = \frac{\bar{\alpha}\, \bar{T}\, g}{\bar{C}_p}
}\end{equation}

\begin{equation}\protect\phantomsection\label{eq:adiabatic-pressure}{
  \frac{d\bar{P}}{dy} = \bar{\rho}\, g
}\end{equation}

where the material properties \(\bar{\rho}\), \(\bar{\alpha}\), and
\(\bar{C}_p\) are determined from the stable equilibrium assemblage
evaluated along the adiabatic reference conditions (Figure S1).

We superimpose Gaussian thermal anomalies of \(\pm\) 500 K on the
adiabatic profile using linear features with 15 km half-width Gaussian
cross-sections and 5 km tanh-tapered ends. Slab anomalies extend 105 km
horizontally and 175 km downward from the top boundary, while plume
anomalies extend 280 km upward from the bottom boundary. We prescribe
constant inflow velocities \(\vec{u}_\mathrm{in}\) of 5 or 10 cm
yr\(^{-1}\) parallel to the thermal anomalies. Boundary conditions on
lateral walls apply zero horizontal velocity with a lithostatic normal
traction, and the (open) top and bottom boundaries apply constant normal
traction equal to the initial lithostatic pressure.

\begin{figure}
\centering
\includegraphics[width=0.45\linewidth,height=\textheight,keepaspectratio,alt={Initial setup for slab (top) and plume (bottom) simulations follows Kerswell et al. (2026), but with an extended 1500 \textbackslash times 1000 km model domain and additional phase transitions (410, 520, and 660; gray bands). Slabs extend 105 km horizontally and 175 km downward from the top boundary, while plumes extend 280 km upward from the bottom boundary. All other boundary conditions are consistent with Kerswell et al. (2026).}]{../figs/setup/initial-setup-diagram.png}
\caption{Initial setup for slab (top) and plume (bottom) simulations
follows Kerswell et al. (2026), but with an extended 1500 \(\times\)
1000 km model domain and additional phase transitions (410, 520, and
660; gray bands). Slabs extend 105 km horizontally and 175 km downward
from the top boundary, while plumes extend 280 km upward from the bottom
boundary. All other boundary conditions are consistent with Kerswell et
al. (2026).}\label{fig:initial-setup-diagram}
\end{figure}

\subsubsection{Material Model}\label{sec:material-model}

\paragraph{Material Properties}\label{sec:material-properties}

We compute material properties with BurnMan with equations of state from
Stixrude \& Lithgow-Bertelloni (2022) for pure Mg ringwoodite,
bridgmanite, and periclase (see
Section~\ref{sec:adiabatic-reference-conditions}), and all properties
except density are held at their adiabatic reference values (Figure S1).
We correct density using a first-order Taylor expansion:

\begin{equation}\protect\phantomsection\label{eq:density-ala}{
  \rho \approx \bar{\rho} \left(1 + \bar{\beta}\, \hat{P} - \bar{\alpha}\, \hat{T} \right)
}\end{equation}

where \(\hat{P} = P - \bar{P}\) and \(\hat{T} = T - \bar{T}\) are the
dynamic pressure and temperature. The volume fraction fields
\(X_\mathrm{wd}\), \(X_\mathrm{ri}\), and \(X_\mathrm{ps}\) account for
phase-transition density jumps and evolve according to the kinetic
models described in Section~\ref{sec:reaction-kinetics}. We hold thermal
conductivity \(\bar{k}\) = 4.0 W m\(^{-1}\)K\(^{-1}\) constant in all
numerical experiments.

\paragraph{Reaction Kinetics}\label{sec:reaction-kinetics}

We treat reaction rates for all three MTZ phase transitions using the
site-saturated nucleation framework of Cahn (1956):

\begin{equation}\protect\phantomsection\label{eq:volume-fraction}{
  X = 1 - \exp\!\left(-N\, \dot{x}\, t \right)
}\end{equation}

\begin{equation}\protect\phantomsection\label{eq:volume-transformation-rate}{
  \dot{X} = N\, \dot{x}\, \left(1 - X \right)
}\end{equation}

where \(X\) is the volume fraction of the product phase, \(\dot{X}\) is
the volume transformation rate, \(N = \text{6.67}\, /\, d\) m\(^{-1}\)
is a geometric factor accounting for grain-boundary nucleation sites
with grain size \(d\), \(\dot{x}\) is the phase-specific growth rate,
and \(t\) is elapsed time after site saturation.

We formulate the 410 and 520 transitions using the interface-controlled
growth model of Hosoya et al. (2005), with macroscopic reaction rates
governed by:

\begin{equation}\protect\phantomsection\label{eq:reaction-rate-410}{
  \dot{X} = Z_\mathrm{ol}\, T\, \exp\!\left(-\frac{H^{\ast} + P V^{\ast}}{R\, T}\right) \left(1 - \exp\!\left[-\frac{\Delta G}{R\, T}\right] \right)\, \left(1 - X \right)
}\end{equation}

where \(H^{\ast}\) = 274 kJ mol\(^{-1}\), \(V^{\ast}\) = 3.0 \(\times\)
10\(^{-6}\) m\(^3\) mol\(^{-1}\), and we hold the kinetic prefactors
fixed at \(Z_\mathrm{ol}\) = \(Z_\mathrm{wd}\) = 1.4 \(\times\) 10\(^4\)
K s\(^{-1}\). These coefficients correspond to relatively fast,
quasi-equilibrium 410 and 520 reaction rates (Kerswell et al., 2026),
and we choose them deliberately to avoid interfering with the kinetic
signal investigated at the 660.

The molar Gibbs free energy difference for each phase transformation is
approximated by:

\begin{equation}\protect\phantomsection\label{eq:gibbs}{
  \Delta G \approx \Delta \bar{G} + \hat{P}\, \Delta \bar{V} - \hat{T}\, \Delta \bar{S}
}\end{equation}

where we compute the thermodynamic quantities \(\Delta \bar{G}\),
\(\Delta \bar{V}\), and \(\Delta \bar{S}\) for each major reaction along
the adiabatic reference profile (Figure S2).

The 660 post-spinel reaction requires a fundamentally different kinetic
formulation. Because one phase breaks down to two phases of different
compositions, diffusion must play a role regardless of whether interface
control has an effect. Microstructural observations show that the
forward ringwoodite \(\Rightarrow\) bridgmanite + periclase
transformation proceeds via diffusion-controlled lamellar decomposition.
Bridgmanite and periclase lamellae nucleate rapidly at grain boundaries
and grow at rates limited by silicon self-diffusion through the parent
ringwoodite (Kubo et al., 2002; Zener, 1946). Following the
diffusion-controlled Zener-Hillert decomposition model (Burke, 1965;
Zener, 1946), the growth rate for the post-spinel assemblage is:

\begin{equation}\protect\phantomsection\label{eq:growth-rate-660}{
  \dot{x} = K^{\prime}\, \exp\!\left(-\frac{E^{\ast}}{R\, T}\right)\, \frac{\Delta G^2}{\bar{V}_\mathrm{ps}^2}
}\end{equation}

where \(K^{\prime}\) is the growth rate prefactor, \(E^{\ast}\) is the
activation energy for Si self-diffusion,
\(\bar{V}_\mathrm{ps}^2 = (V_\mathrm{bg} + V_\mathrm{pe})^2\) is the
squared molar volume of the post-spinel assemblage. We refactor
\(\Delta G^2 = \Delta G \cdot |\Delta G|\) as the thermodynamic driving
force term, which allows for reversible reactions when \(\Delta G\)
\(<\) 0.

Combining \(N\), \(K^{\prime}\), and \(\bar{V}_\mathrm{ps}^2\) into a
single kinetic prefactor
\(Z_\mathrm{ri} = \text{6.67}\, K^{\prime}\, /\, d\, \bar{V}_\mathrm{ps}^2\),
the macroscopic reaction rate for ringwoodite \(\Leftrightarrow\)
bridgmanite + periclase becomes:

\begin{equation}\protect\phantomsection\label{eq:reaction-rate-660}{
  \dot{X} = Z_\mathrm{ri}\, \exp\!\left(-\frac{E^{\ast}}{R\, T}\right)\, \Delta G \cdot |\Delta G| \left(1 - X \right)
}\end{equation}

Equation~\ref{eq:reaction-rate-660} introduces a fundamental shift from
the interface-controlled reaction mechanism at the 410 and 520
(Equation~\ref{eq:reaction-rate-410}) to a diffusion-controlled reaction
mechanism at the 660. Mathematically, the driving force at the 660 is
quadratic (\(\Delta G^2\)), whereas the 410 and 520 follow a linear
approximation near equilibrium \((\Delta G\, /\, RT)\) = 0 before
saturating at high overstepping. Physically, this quadratic dependence
more aggressively suppresses reaction rates near equilibrium for the 660
transition. Conversely, the lack of exponential saturation means the 660
reaction rate remains sensitive to the magnitude of overstepping even
far from equilibrium, unlike the 410 and 520 where the driving force
term trends towards unity. These divergent behaviors underpin the
distinct structural and dynamic signatures discussed in Sections
\ref{sec:results} and \ref{sec:discussion}.

The range of \(Z_\mathrm{ri}\) explored in our experiments {[}1.6
\(\times\) 10\(^{-3}\), 1.6 \(\times\) 10\(^3\){]} mol\(^2\) J\(^{-2}\)
s\(^{-1}\) is calibrated against the high-overpressure experimental
dataset of Kubo et al. (2002) using Monte Carlo regression with a
post-spinel molar volume of \(\bar{V}_\mathrm{ps}\) = 3.381 \(\times\)
10\(^{-5}\) m\(^3\) mol\(^{-1}\) (Figure S3). Our regression yields
\(E^{\ast}\) = 310 \(\pm\) 233 kJ mol\(^{-1}\) and \(\ln(K^{\prime})\) =
-29.8 \(\pm\) 22.1, both consistent with Kubo et al. (2002) within
uncertainty. We therefore adopt \(E^{\ast}\) = 355 kJ mol\(^{-1}\)
(after Kubo et al., 2002), \(\ln(K^{\prime})\) in the range {[}-33.5,
-22.0{]}, and \(d\) in the range {[}1 \(\times\) 10\(^{-3}\), 10
\(\times\) 10\(^{-3}\){]} m (Karato, 1984, 2008), such that
\(Z_\mathrm{ri}\) spans conditions from slow kinetics in coarse-grained
rocks to fast kinetics in fine-grained rocks.

\paragraph{Operator Splitting}\label{sec:operator-splitting}

Since reaction rates for the 410, 520, and 660 transitions can exceed
advective timescales, we employ the first-order operator splitting
scheme to decouple reaction from advection. Volume fractions are updated
in two sequential steps per advection timestep \(\Delta t\): 1) a
reaction substep integrating Equations \ref{eq:reaction-rate-410} and
\ref{eq:reaction-rate-660} over the full \(\Delta t\) using the ARKODE
adaptive Runge-Kutta integrator (Hindmarsh et al., 2005; Reynolds et
al., 2023) with relative tolerance 1 \(\times\) 10\(^{-6}\); and 2) an
advection step:

\begin{equation}\protect\phantomsection\label{eq:composition}{
  \frac{\partial X}{\partial t} + \vec{u} \cdot \nabla X = 0
}\end{equation}

The three kinetically active transitions require simultaneous
integration during the reaction step, with the ARKODE substep size
selected adaptively to satisfy the tolerance across all three reaction
systems.

\subsubsection{Rheological Model}\label{sec:rheological-model}

While the volume fractions \(X_\mathrm{ol}\), \(X_\mathrm{wd}\),
\(X_\mathrm{ri}\), and \(X_\mathrm{ps}\) track the actual mineralogy for
the purpose of density and buoyancy (including metastable phases), we
also define equilibrium fractions \(\phi_\mathrm{ol}\),
\(\phi_\mathrm{wd}\), \(\phi_\mathrm{ri}\), and \(\phi_\mathrm{ps}\).
These fractions represent the thermodynamic state along the reference
adiabat and facilitate smooth transitions in background viscosity. We
define a transition function \(f_j\) for each reaction \(j\):

\begin{equation}\protect\phantomsection\label{eq:transition-function}{
  f_j = \frac{1}{2} \left[ 1 - \tanh \left( \frac{\Delta G_j}{w} \right) \right]
}\end{equation}

where \(w\) = 1000 J mol\(^{-1}\) is a ``Gibbs width'' that smooths the
equilibrium fractions \(\phi\) over a few kilometers, avoiding
discontinuous jumps in background viscosity at sharp equilibrium phase
boundaries. To ensure mass conservation \(\sum \phi_i = \text{1}\), we
partition the stability fractions as follows:

\begin{equation}\protect\phantomsection\label{eq:equilibrium-fractions}{
  \begin{aligned}
    \phi_\mathrm{ol} &= 1 - f_\mathrm{wd} \\
    \phi_\mathrm{wd} &= f_\mathrm{wd}\, (1 - f_\mathrm{ri}) \\
    \phi_\mathrm{ri} &= f_\mathrm{wd}\, f_\mathrm{ri}\, (1 - f_\mathrm{ps}) \\
    \phi_\mathrm{ps} &= f_\mathrm{wd}\, f_\mathrm{ri}\, f_\mathrm{ps}
  \end{aligned}
}\end{equation}

The rheological model incorporates phase-dependent viscosity transitions
as a function of the equilibrium fractions \(\phi\) computed in
Equation~\ref{eq:equilibrium-fractions}. We use a temperature-dependent
Arrhenius viscosity modified by a phase-dependent viscosity jump
prefactor \(b\):

\begin{equation}\protect\phantomsection\label{eq:rheological-model}{
  \eta = b\, \bar{\eta}\, \exp\!\left(-B\,\frac{\hat{T}}{\bar{T}}\right)
}\end{equation}

where \(\bar{\eta}\) = 10\(^{21}\) Pa s is the background viscosity,
\(\bar{T}\) is the reference adiabatic temperature, and
\(\hat{T} = T - \bar{T}\) is the nonadiabatic temperature anomaly. We
vary the activation factor \(B\) between 5 and 10 to control thermal
sensitivity.

We calculate the viscosity jump prefactor \(b\) via a log-linear average
of per-phase viscosity prefactors
(\(\nu_\mathrm{ol}, \nu_\mathrm{wd}, \nu_\mathrm{ri}, \nu_\mathrm{ps}\)),
weighted by the equilibrium fractions \(\phi\) defined in
Equation~\ref{eq:equilibrium-fractions}:

\begin{equation}\protect\phantomsection\label{eq:rheological-jump-factor}{
  \ln(b) = \sum_{i} \phi_i \ln(\nu_i)
}\end{equation}

We adopt this log-linear (geometric) average because it is commonly used
for viscosity contrasts across mineral phase mixtures and avoids the
sensitivity to the single weakest phase that a harmonic (Reuss) average
would introduce (Ji, 2004). By anchoring viscosity jumps to the
equilibrium fractions \(\phi\) rather than the actual (metastable)
mineralogy \(X\), the viscosity structure remains a function of the
adiabatic reference condition. This decoupling prevents spurious
viscosity patches arising from metastable phases within cold slabs,
allowing us to isolate the kinetic contribution to slab impedance via
buoyancy from rheological effects. To maintain numerical stability, we
impose a global viscosity floor and ceiling of 10\(^{19}\) and
10\(^{24}\) Pa s, respectively.

\subsubsection{Numerical Stabilization of Dynamic Pressure
Oscillations}\label{sec:numerical-stabilization-of-dynamic-pressure-oscillations}

Following Gassmöller et al. (2020) and Kerswell et al. (2026), we
exclude the dynamic pressure \(\hat{P}\) contribution to \(\Delta G\) in
the kinetic rate equations (Equations \ref{eq:reaction-rate-410} and
\ref{eq:reaction-rate-660}) while retaining it in the density
formulation (Equation~\ref{eq:density-ala}). In practice, this means
applying \(\Delta G \approx \Delta \bar{G} - \hat{T}\, \Delta \bar{S}\)
to the kinetic rate equations, which preserves the primary thermal
control of the reactions while eliminating pressure-wave feedback. The
quadratic \(\Delta G \cdot |\Delta G|\) driving force amplifies
pressure-wave sensitivity relative to the exponential form used at the
410 and 520, making this stabilization particularly important for the
660. Limitations of this approximation are discussed in
Section~\ref{sec:uncertainties-and-model-limitations}.

\subsubsection{Numerical Experiment Design and
Postprocessing}\label{sec:numerical-experiment-design-and-postprocessing}

We explore the ringwoodite decomposition kinetic parameter space across
eight scenarios spanning two inflow velocities, two rheological strength
contrasts, and two lower-mantle viscosity jumps
(Table~\ref{tbl:scenarios}), running matched slab and plume simulations
across the full tested range of \(Z_\mathrm{ri}\) within every scenario.

We characterize 660 structure using two diagnostic measures extracted
along a representative depth profile through each simulation.
`Displacement' is the difference between the depth at which the
post-spinel volume fraction \(X_\mathrm{ps}\) reaches 0.9 and the
nominal equilibrium post-spinel depth. `Width' is the depth interval
between \(X_\mathrm{ps}\) = 0.9 and \(X_\mathrm{ps}\) = 0.1. We also
report the maximum post-spinel reaction rate \(\dot{X}_\mathrm{max}\)
and the maximum descent (or ascent) velocity \(\vec{v}_\mathrm{max}\)
along the same profile, extracted within the phase transition zone
(defined by the measured 660 width). These later two quantities allow us
to define a characteristic reaction completion length scale (in km):

\begin{equation}\protect\phantomsection\label{eq:characteristic-length-scale}{
  L = 10\, \frac{\vec{v}_\mathrm{max}}{\dot{X}_\mathrm{max}}
}\end{equation}

which approximates the depth traversed during one characteristic
reaction timescale \(\tau = 1/\dot{X}_\mathrm{max}\) (Ma). \(L\) and the
measured 660 width are related but not identical. \(L\) is a local
estimate based on the maximum reaction rate and maximum vertical
velocity through the 660, whereas width measures the full spatial extent
of the transition between \(X_\mathrm{ps}\) = 0.1 and 0.9. The two
metrics diverge where reaction rate varies strongly within the phase
transition zone.

In slab simulations, \(L\) and \(\vec{v}_\mathrm{max}\) are subsequently
used to heuristically classify each model into three regimes using fixed
threshold values. The quasi-equilibrium regime corresponds to \(L\)
\(\leq\) 5 km and \(\vec{v}_\mathrm{max}\) \(>\) 0.5 cm yr\(^{-1}\); the
intermediate regime corresponds to \(L\) \(>\) 5 km and
\(\vec{v}_\mathrm{max}\) \(>\) 0.5 cm yr\(^{-1}\); and the stagnation
regime correspond to \(\vec{v}_\mathrm{max}\) \(\leq\) 0.5 cm
yr\(^{-1}\).

\begin{longtable}[]{@{}lrrr@{}}
\caption{Rheological and kinematic parameters for the eight tested
scenarios. Units are \(\vec{u}_\mathrm{in}\): cm yr\(^{-1}\), \(B\):
dimensionless, \(b\):
dimensionless.}\label{tbl:scenarios}\tabularnewline
\toprule\noalign{}
Scenario & \(\vec{u}_\mathrm{in}\) & \(B\) & \(b\) \\
\midrule\noalign{}
\endfirsthead
\toprule\noalign{}
Scenario & \(\vec{u}_\mathrm{in}\) & \(B\) & \(b\) \\
\midrule\noalign{}
\endhead
\bottomrule\noalign{}
\endlastfoot
0 & 5 & 5 & 50 \\
1 & 5 & 10 & 50 \\
2 & 10 & 5 & 50 \\
3 & 10 & 10 & 50 \\
4 & 5 & 5 & 1 \\
5 & 5 & 10 & 1 \\
6 & 10 & 5 & 1 \\
7 & 10 & 10 & 1 \\
\end{longtable}

\section{Results}\label{sec:results}

\subsection{Simulation Snapshots: Slabs and
Plumes}\label{sec:simulation-snapshots}

Figures \ref{fig:slab-composition-660-b01} and
\ref{fig:plume-composition-660-b01} show slab and plume simulations for
scenario 4 (Table~\ref{tbl:scenarios}) after 100 Ma and 50 Ma,
respectively. This reference scenario illustrates the range of 660
structures produced across three kinetic conditions characterizing slow
to fast reaction rates. Ringwoodite decomposition kinetics strongly
control slab dynamics while affecting plume structure more modestly,
consistent with the asymmetric pattern documented at the 410 by Kerswell
et al. (2026).

In slab simulations, three regimes emerge with decreasing
\(Z_\mathrm{ri}\). In the quasi-equilibrium regime
(Figure~\ref{fig:slab-composition-660-b01}: bottom row;
\(Z_\mathrm{ri}\) = 1.6 \(\times\) 10\(^3\) mol\(^2\) J\(^{-2}\)
s\(^{-1}\)), the slab descends through the 660 at 1.6 cm yr\(^{-1}\) and
the 660 is narrow (6 km) and modestly displaced (25 km), consistent with
Clapeyron-slope predictions for the imposed thermal anomaly. In the
intermediate regime (Figure~\ref{fig:slab-composition-660-b01}: middle
row; \(Z_\mathrm{ri}\) = 6.0 \(\times\) 10\(^{-1}\) mol\(^2\) J\(^{-2}\)
s\(^{-1}\)), metastable ringwoodite accumulates and reduces descent
through the 660 to 1.1 cm yr\(^{-1}\), deepening the 660 by 40 km and
widening it to 12 km. In the stagnation regime
(Figure~\ref{fig:slab-composition-660-b01}: top row; \(Z_\mathrm{ri}\) =
1.6 \(\times\) 10\(^{-3}\) mol\(^2\) J\(^{-2}\) s\(^{-1}\)), persistent
metastable buoyancy arrests descent (\(\vec{v}_\mathrm{max}\) = 0.25 cm
yr\(^{-1}\)) and the slab ponds above the 660, displacing it by 48 km
and broadening it to 11 km.

In plume simulations, the 660 remains uplifted and relatively sharp
across all tested kinetic conditions. Under the slowest reaction rates
(Figure~\ref{fig:plume-composition-660-b01}: top row; \(Z_\mathrm{ri}\)
= 1.6 \(\times\) 10\(^{-3}\) mol\(^2\) J\(^{-2}\) s\(^{-1}\)), the plume
ascends through the 660 at 3 cm yr\(^{-1}\) with the 660 displaced by 14
km upward and widening to 9 km. Under intermediate and fast kinetic
conditions (Figure~\ref{fig:plume-composition-660-b01}: middle and
bottom rows; \(Z_\mathrm{ri}\) = 6.0 \(\times\) 10\(^{-1}\) and 1.6
\(\times\) 10\(^3\) mol\(^2\) J\(^{-2}\) s\(^{-1}\)), the 660 narrows to
7 km and is uplifted 10 km, with ascent velocities through the 660 of 4
cm yr\(^{-1}\). Elevated plume temperatures maintain
\(\dot{X}_\mathrm{max}\) above 3 Ma\(^{-1}\) even at the slowest tested
\(Z_\mathrm{ri}\), suppressing extended metastable zones like those
formed in cold slabs. This kinetic insensitivity of plume dynamics and
660 structures holds qualitatively across the full tested parameter
space and contrasts sharply with the strong kinetic control on slab
behavior.

\begin{figure}
\centering
\includegraphics[width=1\linewidth,height=\textheight,keepaspectratio,alt={Reference slab simulations with moderate inflow velocity, intermediate slab strength, and no lower-mantle viscosity jump (scenario 4; Table~). Rows show stagnation (top row: Z\_\textbackslash mathrm\{ri\} = 1.6 \textbackslash times 10\^{}\{-3\} mol\^{}2 J\^{}\{-2\} s\^{}\{-1\}), intermediate (middle row: Z\_\textbackslash mathrm\{ri\} = 6.0 \textbackslash times 10\^{}\{-1\} mol\^{}2 J\^{}\{-2\} s\^{}\{-1\}), and quasi-equilibrium (bottom row: Z\_\textbackslash mathrm\{ri\} = 1.6 \textbackslash times 10\^{}3 mol\^{}2 J\^{}\{-2\} s\^{}\{-1\}) kinetic regimes after 100 Ma. Columns show dynamic temperature \textbackslash hat\{T\} (left column), dynamic density \textbackslash hat\{\textbackslash rho\} (middle column), and vertical velocity \textbackslash vec\{u\}\_y (right column). Thick black lines (left column) highlight the 10\% and 90\% wadsleyite and post-spinel volume fraction contours used to define 660 displacement and width. The background grid pattern indicates ASPECT's adaptive mesh refinement near high thermal gradients and phase transitions. White trace (left column) indicates the representative depth profile used to extract data, and the white arrows indicate where the 660 structure is evaluated.}]{../figs/simulation/compositions/slab-set1-composition-b01-0100.png}
\caption{Reference slab simulations with moderate inflow velocity,
intermediate slab strength, and no lower-mantle viscosity jump (scenario
4; Table~\ref{tbl:scenarios}). Rows show stagnation (top row:
\(Z_\mathrm{ri}\) = 1.6 \(\times\) 10\(^{-3}\) mol\(^2\) J\(^{-2}\)
s\(^{-1}\)), intermediate (middle row: \(Z_\mathrm{ri}\) = 6.0
\(\times\) 10\(^{-1}\) mol\(^2\) J\(^{-2}\) s\(^{-1}\)), and
quasi-equilibrium (bottom row: \(Z_\mathrm{ri}\) = 1.6 \(\times\)
10\(^3\) mol\(^2\) J\(^{-2}\) s\(^{-1}\)) kinetic regimes after 100 Ma.
Columns show dynamic temperature \(\hat{T}\) (left column), dynamic
density \(\hat{\rho}\) (middle column), and vertical velocity
\(\vec{u}_y\) (right column). Thick black lines (left column) highlight
the 10\% and 90\% wadsleyite and post-spinel volume fraction contours
used to define 660 displacement and width. The background grid pattern
indicates ASPECT's adaptive mesh refinement near high thermal gradients
and phase transitions. White trace (left column) indicates the
representative depth profile used to extract data, and the white arrows
indicate where the 660 structure is
evaluated.}\label{fig:slab-composition-660-b01}
\end{figure}

\begin{figure}
\centering
\includegraphics[width=1\linewidth,height=\textheight,keepaspectratio,alt={Reference plume simulations with moderate inflow velocity, intermediate plume strength, and no lower-mantle viscosity jump (scenario 4; Table~). Rows show ultra-sluggish (top row: Z\_\textbackslash mathrm\{ri\} = 1.6 \textbackslash times 10\^{}\{-3\} mol\^{}2 J\^{}\{-2\} s\^{}\{-1\}), intermediate (middle row: Z\_\textbackslash mathrm\{ri\} = 6.0 \textbackslash times 10\^{}\{-1\} mol\^{}2 J\^{}\{-2\} s\^{}\{-1\}), and fast (bottom row: Z\_\textbackslash mathrm\{ri\} = 1.6 \textbackslash times 10\^{}3 mol\^{}2 J\^{}\{-2\} s\^{}\{-1\}) kinetic conditions after 50 Ma. All other annotations follow Figure~. The 660 is uplifted across all kinetic conditions, consistent with the negative Clapeyron slope, but 660 width increases and plume ascent slows measurably under the slowest kinetic conditions (top row).}]{../figs/simulation/compositions/plume-set1-composition-b01-0050.png}
\caption{Reference plume simulations with moderate inflow velocity,
intermediate plume strength, and no lower-mantle viscosity jump
(scenario 4; Table~\ref{tbl:scenarios}). Rows show ultra-sluggish (top
row: \(Z_\mathrm{ri}\) = 1.6 \(\times\) 10\(^{-3}\) mol\(^2\) J\(^{-2}\)
s\(^{-1}\)), intermediate (middle row: \(Z_\mathrm{ri}\) = 6.0
\(\times\) 10\(^{-1}\) mol\(^2\) J\(^{-2}\) s\(^{-1}\)), and fast
(bottom row: \(Z_\mathrm{ri}\) = 1.6 \(\times\) 10\(^3\) mol\(^2\)
J\(^{-2}\) s\(^{-1}\)) kinetic conditions after 50 Ma. All other
annotations follow Figure~\ref{fig:slab-composition-660-b01}. The 660 is
uplifted across all kinetic conditions, consistent with the negative
Clapeyron slope, but 660 width increases and plume ascent slows
measurably under the slowest kinetic conditions (top
row).}\label{fig:plume-composition-660-b01}
\end{figure}

\subsection{Structure of the 660: Displacement and
Width}\label{sec:structure-of-the-660}

Figure~\ref{fig:660-structure} summarizes 660 displacement and width as
a function of maximum reaction rate \(\dot{X}_\mathrm{max}\) across the
full suite of plume and slab simulations, and
Figure~\ref{fig:660-structure-comp} maps all results across the full
kinetic-rheological-kinematic parameter space (Table S1). Together, the
figures reveal three kinetic regimes in slabs, near-invariant 660
structure in plumes, plus independent and coupled
kinetic-kinematic-rheological effects on 660 structure and slab descent.

Plume 660 structure remains largely independent of kinetics across the
full tested range of kinetic prefactors \(Z_\mathrm{ri}\)
(Figure~\ref{fig:660-structure}). Upward displacements vary only from
10--22 km and widths from 5--10 km across six orders of magnitude
variation in \(Z_\mathrm{ri}\), and ascent velocities through the 660
remain within 2--8 cm yr\(^{-1}\) regardless of kinetics. These are
relatively narrow ranges compared to the order-of-magnitude changes seen
in slabs. Even at the slowest tested conditions (\(Z_\mathrm{ri}\) = 1.6
\(\times\) 10\(^{-3}\) mol\(^2\) J\(^{-2}\) s\(^{-1}\)), maximum
reaction rates in plumes reach 3--7 Ma\(^{-1}\) (\(\tau \lesssim 0.3\)
Ma; \(L \lesssim 10\) km), maintaining a narrow (7--10 km) and
consistently uplifted (14--22 km) 660. Plume 660 widths are
systematically broader than plume 410 widths at equivalent kinetic
conditions, but only marginally (cf.~circle and triangle datapoints in
Figure~\ref{fig:660-structure}).

Slab 660 structure changes substantially across three kinetic regimes
(Figure~\ref{fig:660-structure-comp}). Regime boundaries follow two
heuristic thresholds: 1) \(L\) \(\lesssim\) 5 km separates
quasi-equilibrium from intermediate regimes, and 2)
\(\vec{v}_\mathrm{max}\) \(\lesssim\) 0.5 cm yr\(^{-1}\) separates
intermediate from stagnation regimes (see
Section~\ref{sec:numerical-experiment-design-and-postprocessing}). In
the quasi-equilibrium regime (region 1 in
Figure~\ref{fig:660-structure-comp}), \(\tau\) remains below 2 Ma and
slabs descend through the 660 at 0.5--3.2 cm yr\(^{-1}\), with the 660
displaced 22--36 km downward and 4--12 km wide. In the intermediate
regime (region 2 in Figure~\ref{fig:660-structure-comp}), \(\tau\) spans
0.5--13.1 Ma and descent velocities through the 660 of 0.4--2.4 cm
yr\(^{-1}\) produce continuous but progressively slower slab
penetration, with the 660 deepening by 40--80 km and widening to 9--34
km. In the stagnation regime (region 3 in
Figure~\ref{fig:660-structure-comp}), descent velocities through the 660
of 0.1--0.5 cm yr\(^{-1}\) and \(\tau\) up to 26 Ma leave the 660
displaced 27--80 km downward with widths of 6--34 km. In contrast to the
abrupt re-sharpening threshold at the 410 (Kerswell et al., 2026), 660
displacement and width increase monotonically with decreasing
\(Z_\mathrm{ri}\) across all three regimes, with no re-sharpening at any
tested kinetic prefactor (Figure~\ref{fig:660-structure}). The direction
of 660 displacement also differs from the 410. At the 410, displacement
changes sign from uplifted under quasi-equilibrium conditions to
depressed in the ultra-sluggish regime (Kerswell et al., 2026), whereas
at the 660, slower kinetics amplify downward displacement monotonically
across all three regimes (cf.~circle and triangle datapoints in
Figure~\ref{fig:660-structure}).

\begin{figure}
\centering
\includegraphics[width=0.7\linewidth,height=\textheight,keepaspectratio,alt={Measured 660 width (top) and displacement (bottom) versus maximum reaction rates in plume (left) and slab (right) simulations after 50 Ma and 100 Ma, respectively. Triangles show the 660 results from this study (Table S1) while circles show the 410 results from Kerswell et al. (2026) for comparison. The 660 structure near plumes shows consistent upward displacements and moderate widths with minimal dependence on \textbackslash dot\{X\}\_\textbackslash mathrm\{max\} across six orders of magnitude variation. Unlike the 410, the 660 structure near slabs (right column) changes monotonically with decreasing \textbackslash dot\{X\}\_\textbackslash mathrm\{max\} and no re-sharpening occurs at ultra-sluggish kinetic conditions. The modest scatter of the 660 results (triangles) at fixed \textbackslash dot\{X\}\_\textbackslash mathrm\{max\} reflects the additional, systematically varying influence of inflow velocity (\textbackslash vec\{u\}\_\textbackslash mathrm\{in\}), rheological contrast (B), and lower-mantle viscosity jump (b). These effects are explored more completely in Figure~.}]{../figs/structure/composition.png}
\caption{Measured 660 width (top) and displacement (bottom) versus
maximum reaction rates in plume (left) and slab (right) simulations
after 50 Ma and 100 Ma, respectively. Triangles show the 660 results
from this study (Table S1) while circles show the 410 results from
Kerswell et al. (2026) for comparison. The 660 structure near plumes
shows consistent upward displacements and moderate widths with minimal
dependence on \(\dot{X}_\mathrm{max}\) across six orders of magnitude
variation. Unlike the 410, the 660 structure near slabs (right column)
changes monotonically with decreasing \(\dot{X}_\mathrm{max}\) and no
re-sharpening occurs at ultra-sluggish kinetic conditions. The modest
scatter of the 660 results (triangles) at fixed \(\dot{X}_\mathrm{max}\)
reflects the additional, systematically varying influence of inflow
velocity (\(\vec{u}_\mathrm{in}\)), rheological contrast (\(B\)), and
lower-mantle viscosity jump (\(b\)). These effects are explored more
completely in
Figure~\ref{fig:660-structure-comp}.}\label{fig:660-structure}
\end{figure}

Figure~\ref{fig:660-structure-comp} further isolates four trends in slab
simulations across the eight tested scenarios
(Table~\ref{tbl:scenarios}). The first is an independent effect, while
the latter three reflect coupled interactions between two or more
parameters.

The kinetic prefactor \(Z_\mathrm{ri}\) controls maximum reaction rate
independently of rheology or kinematics. Maximum reaction rate decreases
monotonically and proportionally with \(Z_\mathrm{ri}\) across all eight
scenarios, spanning approximately 0.04--135 Ma\(^{-1}\)
(Figure~\ref{fig:660-structure-comp}), consistent with the direct
\(Z_\mathrm{ri}\) scaling in Equation~\ref{eq:reaction-rate-660}.

The lower-mantle viscosity jump (\(b\) = 50, compared to no jump at
\(b\) = 1) exerts a first-order control on slab descent velocity through
the 660 (see Figures S4--S7), shifting the stagnation regime boundary
toward higher \(Z_\mathrm{ri}\). Comparing scenarios with \(b\) = 50
(scenarios 0--3) to \(b\) = 1 (scenarios 4--7) at equivalent \(B\) and
\(\vec{u}_\mathrm{in}\), the stagnation regime boundary shifts by up to
four orders of magnitude in \(Z_\mathrm{ri}\) for the weakest and
slowest slab configurations (cf.~scenarios 0 and 4 in
Figure~\ref{fig:660-structure-comp}). This effect is most pronounced for
weaker slabs (\(B\) = 5). For instance, scenario 0
(\(\vec{u}_\mathrm{in}\) = 5 cm yr\(^{-1}\), \(B\) = 5, \(b\) = 50)
falls at or below the 0.5 cm yr\(^{-1}\) stagnation threshold for all
tested \(Z_\mathrm{ri}\), effectively stagnating across the entire
kinetic range. For stronger slabs (\(B\) = 10; scenarios 1 and 3), the
effect of a lower-mantle viscosity jump on the stagnation regime
boundary is substantially reduced. Thus, while the presence of a
viscosity jump consistently promotes stagnation, its absolute impact is
modulated by an interaction between mantle and slab strength contrasts
(\(b\) versus \(B\)).

Inflow velocity and kinetics also interact to influence descent velocity
through the 660 and 660 structure. Higher inflow velocity
(\(\vec{u}_\mathrm{in}\) = 10 cm yr\(^{-1}\)) shifts the stagnation
regime boundary toward lower \(Z_\mathrm{ri}\) relative to
\(\vec{u}_\mathrm{in}\) = 5 cm yr\(^{-1}\) at equivalent \(B\) and
\(b\). This happens most clearly for stronger slabs (cf.~scenarios 1 and
3; 5 and 7 in Figure~\ref{fig:660-structure-comp}), where the shift
reaches 1--2 orders of magnitude. At the same time, faster inflow
produces larger displacements and widths across all regimes;
\(\vec{u}_\mathrm{in}\) = 10 cm yr\(^{-1}\) scenarios (2, 3, 6, 7) reach
22--80 km displacement and 4--34 km width, compared to 22--59 km and
4--21 km for \(\vec{u}_\mathrm{in}\) = 5 cm yr\(^{-1}\) scenarios (0, 1,
4, 5). Both the stagnation regime boundary shift and the increase in 660
displacement and width reflect that faster descent advects more
metastable ringwoodite to greater depths before ponding.

In parallel with the coupled interactions above, slab strength and
descent velocity through the 660 operate in tandem to further modulate
660 structure. Weaker slabs (\(B\) = 5; scenarios 0, 2, 4, 6) decelerate
more readily than stronger slabs (\(B\) = 10; scenarios 1, 3, 5, 7),
stagnating at up to two orders of magnitude higher \(Z_\mathrm{ri}\) and
shifting both regime boundaries toward faster kinetic conditions
(Figure~\ref{fig:660-structure-comp}). In the intermediate regime,
stronger slabs produce substantially larger 660 displacements (40--80
km; mean 57 km) than weaker slabs (40--56 km; mean 47 km) at equivalent
\(Z_\mathrm{ri}\). The fastest and strongest slabs (scenarios 3 and 7)
resist stagnation, even under ultra-sluggish kinetic conditions
(\(Z_\mathrm{ri}\) \(\lesssim\) 1.2 \(\times\) 10\(^{-2}\) mol\(^2\)
J\(^{-2}\) s\(^{-1}\)), producing the largest displacements (78--80 km)
and widths (26--34 km) in our simulations. At the same
\(Z_\mathrm{ri}\), weaker counterparts (scenarios 2 and 6) are already
stagnant, with notably smaller displacements (68--75 km). Taken
together, weaker slabs are more prone to stagnation but produce smaller
maximum displacements, while stronger slabs resist stagnation longer and
ultimately generate the widest and most deeply displaced 660s when their
descent through the 660 slows---a trend that reflects coupled kinetic,
kinematic, and rheological controls on 660 structure.

\begin{figure}
\centering
\includegraphics[width=0.7\linewidth,height=\textheight,keepaspectratio,alt={Variation in 660 structure and slab descent rate through the 660 across kinetic (Z\_\textbackslash mathrm\{ri\}), kinematic (\textbackslash vec\{u\}\_\textbackslash mathrm\{in\}), and rheological (B, b) parameter space. Panels show maximum reaction rate (top left), maximum vertical velocity through the 660 (top right), 660 width (bottom left), and 660 displacement (bottom right) as functions of the kinetic prefactor Z\_\textbackslash mathrm\{ri\} (horizontal axis, log scale) and remaining parameter space reduced to eight scenarios {[}\textbackslash vec\{u\}\_\textbackslash mathrm\{in\}, B, b{]}: 0) {[}5, 5, 50{]}; 1) {[}5, 10, 50{]}; 2) {[}10, 5, 50{]}; 3) {[}10, 10, 50{]}; 4) {[}5, 5, 1{]}; 5) {[}5, 10, 1{]}; 6) {[}10, 5, 1{]}; 7) {[}10, 10, 1{]} (Table~). Units are \textbackslash vec\{u\}\_\textbackslash mathrm\{in\}: cm yr\^{}\{-1\}, B: dimensionless, b: dimensionless. Each tile represents the measured value after 100 Ma (Table S1). Black and white lines delineate transitions between regime behaviors, and their positions shift with B, b, and \textbackslash vec\{u\}\_\textbackslash mathrm\{in\}. Regime labels correspond to: (1) quasi-equilibrium, (2) intermediate, and (3) stagnation.}]{../figs/structure/tiles-660.png}
\caption{Variation in 660 structure and slab descent rate through the
660 across kinetic (\(Z_\mathrm{ri}\)), kinematic
(\(\vec{u}_\mathrm{in}\)), and rheological (\(B\), \(b\)) parameter
space. Panels show maximum reaction rate (top left), maximum vertical
velocity through the 660 (top right), 660 width (bottom left), and 660
displacement (bottom right) as functions of the kinetic prefactor
\(Z_\mathrm{ri}\) (horizontal axis, log scale) and remaining parameter
space reduced to eight scenarios {[}\(\vec{u}_\mathrm{in}\), \(B\),
\(b\){]}: \textbf{0)} {[}5, 5, 50{]}; \textbf{1)} {[}5, 10, 50{]};
\textbf{2)} {[}10, 5, 50{]}; \textbf{3)} {[}10, 10, 50{]}; \textbf{4)}
{[}5, 5, 1{]}; \textbf{5)} {[}5, 10, 1{]}; \textbf{6)} {[}10, 5, 1{]};
\textbf{7)} {[}10, 10, 1{]} (Table~\ref{tbl:scenarios}). Units are
\(\vec{u}_\mathrm{in}\): cm yr\(^{-1}\), \(B\): dimensionless, \(b\):
dimensionless. Each tile represents the measured value after 100 Ma
(Table S1). Black and white lines delineate transitions between regime
behaviors, and their positions shift with \(B\), \(b\), and
\(\vec{u}_\mathrm{in}\). Regime labels correspond to: (1)
quasi-equilibrium, (2) intermediate, and (3)
stagnation.}\label{fig:660-structure-comp}
\end{figure}

\section{Discussion}\label{sec:discussion}

\subsection{Uncertainties and Model
Limitations}\label{sec:uncertainties-and-model-limitations}

The primary quantitative uncertainty for the post-spinel reaction stems
from the activation energy \(E^{\ast}\) for Si self-diffusion and the
kinetic prefactor \(Z_\mathrm{ri}\), which bundles the growth rate
prefactor \(K^{\prime}\) and grain size \(d\). Monte Carlo regression of
the Kubo et al. (2002) experimental dataset yields a mean activation
energy of 310 \(\pm\) 233 kJ mol\(^{-1}\) and a mean \(\ln(K^{\prime})\)
of -29.8 \(\pm\) 22.1 (Figure S3). These uncertainties span several
orders of magnitude, reflecting incomplete knowledge of Si
self-diffusion in ringwoodite at natural PT conditions, dependence on
grain size and deformation history, and poorly characterized water
content effects on post-spinel diffusion rates. Our simulations
therefore bracket plausible kinetic conditions rather than targeting
specific natural scenarios, and all three kinetic regimes identified for
slabs fall within this uncertainty envelope.

Another key driver of uncertainty is that the kinetic experiments of
Kubo et al. (2002) were performed at highly-overstepped metastable
conditions (\(>\) 1 GPa), requiring significant extrapolation to
near-equilibrium conditions in our simulations. The quadratic dependence
of ringwoodite decomposition on excess Gibbs energy
(Equation~\ref{eq:reaction-rate-660}) results in vanishingly small
reaction rates at small pressure oversteps in our simulations. However,
it is possible that an alternative, faster transformation mechanism
occurs in natural systems at near-equilibrium conditions. This requires
further kinetic experiments at small pressure oversteps to determine.

A related simplification concerns the reaction's directionality. Our
bidirectional implementation of the reverse reaction (ringwoodite
regrowth from bridgmanite + periclase) treats the forward and reverse
transformations symmetrically, which represents a simplification.
Experiments document that the reverse reaction proceeds substantially
more slowly than the forward decomposition, particularly below 2000 K
(Lessing et al., 2022), introducing a kinetic asymmetry not captured by
our symmetric \(\Delta G \cdot |\Delta G|\) formulation in
Equation~\ref{eq:reaction-rate-660}. Our simplification likely
overestimates reverse reaction rates in ascending plumes and
underestimates metastable bridgmanite + periclase persistence in
upwelling material. The slight broadening of plume 660 widths at the
slowest kinetics in our simulations (up to 10 km) may partly reflect
this effect, but a rigorous treatment requires a separate kinetic rate
equation for the reverse reaction.

This kinetic framework is itself coupled to thermodynamics in ways we
only partially resolve. While our material model implicitly includes
latent heating near equilibrium, the latent-heating effect at
overstepped metastable conditions merits particular caution for the
post-spinel reaction. Because the ringwoodite \(\Leftrightarrow\)
bridgmanite + periclase Clapeyron slope is negative, the (endothermic)
latent-heat effect close to equilibrium reinforces kinetic inhibition.
At highly-overstepped metastable conditions, however, the free energy
from mechanical work \(P\Delta V\) (rather than \(T\Delta S\)) becomes
operative and this reinforcement is lost. This contrasts with the 410,
where both the (exothermic) latent-heat effect and \(P\Delta V\)
contributions promote faster kinetics. Fully resolving this feedback
would likely strengthen the regime contrasts we report here and merits
dedicated investigation. We also define Gibbs free energy based on the
reference pressure (Equation~\ref{eq:gibbs}), neglecting deviatoric
stresses at individual grain interfaces that directly influence reaction
pathways and chemical potentials (Wheeler, 2015a, 2020). While the
reference pressure can provide a reasonable approximation under certain
conditions (Wheeler, 2015b), an improved approach would be to use the
dynamic pressure, i.e.~the mean normal stress actually experienced by
the rock. A rigorous treatment will require consideration of more
localized stress-dependent chemical potentials (work in progress).

Beyond these kinetic and thermodynamic uncertainties, our compositional
choices introduce additional simplifications. Our pure Mg end-member
composition neglects Fe-Mg solid solutions, akimotoite, garnet, and
other transition-zone phases present in natural systems (Ita \&
Stixrude, 1992; Papanagnou et al., 2023; Xu et al., 2008). Iron is
unlikely to alter our conclusions substantially because Si
self-diffusion limits ringwoodite decomposition. Iron content should
therefore have only a minor effect on reaction kinetics, and the narrow
(\(<\) 1 kbar) Fe-Mg post-spinel phase loop ({Ishii et al.}, 2019) means
metastable Fe-bearing ringwoodite likely transforms directly to the
stable bridgmanite + periclase assemblage, keeping the transition
effectively univariant as in our pure Mg model (Ming et al., 2006). In
the coldest slab interiors, the akimotoite \(\Leftrightarrow\)
bridgmanite transition may dominate over the post-spinel transition
({Chanyshev et al.}, 2022; Cottaar \& Deuss, 2016), producing additional
deepening beyond our kinetic model. This reaction is nevertheless
expected to be comparatively rapid because the ilmenite and perovskite
structures are crystallographically closely related (Navrotsky, 1998),
permitting a diffusionless or martensite-like reaction mechanism.

A final set of simplifications concerns the large-scale viscosity
structure and dimensionality of our simulations. Our simulations test
lower-mantle viscosity jumps of \(b\) = 1 and \(b\) = 50 to partially
characterize how viscosity structure interacts with kinetic inhibition
(Section~\ref{sec:structure-of-the-660}; see Figures S4--S7). This
approach does not span the full plausible range of lower-mantle
viscosity contrasts (estimated at roughly 10--100\(\times\), Karato,
2008; {Goes et al.}, 2017), and quantitative predictions remain
model-dependent. More broadly, natural subduction involves the
simultaneous action of viscosity structure, the Clapeyron slope, plate
forcing, and kinetics (Garel et al., 2014; {Goes et al.}, 2017), and
isolating each contribution requires independent observational
constraints. Grain-size distribution at the post-spinel reaction front
(Kubo et al., 2000; Mao \& Zhong, 2021) could also generate viscosity
contrasts that augments stagnation beyond what our rheological model
captures. Our simple temperature-dependent rheological model omits
grain-size evolution that affect both flow dynamics and reaction
kinetics (Hirth \& Kohlstedf, 2003; Karato, 1984; Yamazaki et al.,
2005). Finally, our 2D simulations neglect 3D slab geometry, trench
migration, slab rollback, and lateral flow diversion. In 3D, a slab that
retains its dip and is carried laterally away from the 660 by rollback
could sustain faster descent than the vertically ponding slabs in our 2D
domain, since less metastable material would accumulate locally beneath
any one point on the slab. These 3D effects, alongside slab morphology
and residence time more broadly, warrant further investigation (Garel et
al., 2014; {Sime et al.}, 2024). Despite these limitations, our results
capture the first-order kinetic effects governing 660 structure and slab
dynamics across three distinct regimes.

\subsection{Contrasting Reaction Mechanisms at the 410 and
660}\label{sec:contrasting-reaction-mechanisms-at-the-410-and-660}

The three kinetic regimes identified for the 660 differ qualitatively
from those at the 410, reflecting the different reaction mechanisms that
govern each transition (Equations \ref{eq:volume-fraction} to
\ref{eq:reaction-rate-660}). At the 410, the interface-controlled
driving-force term \((1 - \exp[-\Delta G / RT])\) saturates toward unity
once \(|\Delta G| \gg RT\) (Equation~\ref{eq:reaction-rate-410}). Beyond
this saturation point, only the Arrhenius term
\(\exp[-(H^\ast + PV^\ast)/RT]\) controls reaction progress. In a cold,
stagnated slab, this factor remains negligible until temperatures rise
enough to activate transformation, at which point the reaction completes
abruptly over a narrow depth interval, re-sharpening the 410. This
threshold produced the three regimes (quasi-equilibrium, intermediate,
and ultra-sluggish with re-sharpening) with distinct 410 seismic
signatures (Kerswell et al., 2026).

At the 660, the diffusion-controlled driving-force term
\(\Delta G \cdot |\Delta G|\) grows quadratically without bound as
ringwoodite is carried below its stability field. Larger overstepping
always drives larger (silicon) chemical potential gradients across the
decomposition lamellae and thus faster transformation (Kubo et al.,
2002; Zener, 1946). No saturation ceiling exists, so the driving force
never yields to the Arrhenius term alone. As a result,
\(\dot{X}_\mathrm{max}\) increases progressively with depth and the 660
widens and deepens continuously across all three regimes without
triggering an abrupt transformation. The distinction between the
intermediate and stagnation regimes at the 660 is therefore a matter of
degree, marked by progressively longer \(\tau\), larger \(L\), slower
descent, and broader, more deeply displaced discontinuities, rather than
the abrupt threshold that separates the ultra-sluggish regime at the 410
(Figure~\ref{fig:660-structure}).

This difference in reaction mechanism also determines long-term slab
behavior. Metastable olivine accumulated in the ultra-sluggish 410
regime eventually transforms in a rapid pulse that can disrupt transient
stalling or ponding (Kerswell et al., 2026). Metastable ringwoodite in
the 660 stagnation regime is consumed gradually (\(\tau\) up to 26 Ma),
sustaining persistent buoyancy that resists penetration over much longer
timescales. Kinetically stagnated slabs at the 660 are therefore more
stable than those at the 410, consistent with the long-lived
sub-horizontal ponding observed in the western Pacific (Fukao \&
Obayashi, 2013; Tauzin et al., 2017).

Near equilibrium, the two reaction mechanisms also differ. The quadratic
driving force at the 660 strongly suppresses reaction rates near the
phase boundary because \(\Delta G \cdot |\Delta G|\) approaches zero as
\(\Delta G^2\). The linear driving force term governing the 410
\((1 - \exp[-\Delta G / RT]) \approx \Delta G / RT\) suppresses reaction
rates less strongly when \(\Delta G\) approaches zero. This inherent
near-equilibrium suppression broadens the 660 transition zone slightly
relative to the 410, even under quasi-equilibrium conditions (cf.~circle
and triangle datapoints in Figure~\ref{fig:660-structure}).

\subsection{Kinetic Amplification of Temperature-Driven 660
Displacement}\label{sec:kinetic-amplification-of-temperature-driven-660-displacement}

Within the pure Mg, garnet-free system we model, equilibrium
thermodynamics predicts downward 660 displacement from two principal
sources: 1) the negative Clapeyron slope shifting the post-spinel
reaction to higher pressure beneath cold slabs, and 2) the akimotoite
\(\Leftrightarrow\) bridgmanite transition generating additional
depression of 30--90 km in the coldest slab interiors ({Chanyshev et
al.}, 2022; Cottaar \& Deuss, 2016). Our quasi-equilibrium slab
simulations recover this thermal signal (Figures \ref{fig:660-structure}
and \ref{fig:660-structure-comp}), showing 22--36 km displacements
consistent with Clapeyron-slope predictions for a -500 K anomaly (Bina
\& Helffrich, 1994; Flanagan \& Shearer, 1998). Kinetics then amplify
this baseline progressively (Figure~\ref{fig:seismic-profiles}), as the
deepest stagnation regime displacements (up to 80 km) reach roughly 2 to
4 times the shallowest quasi-equilibrium displacement (22 km), depending
on rheological and kinematic conditions (Table S1). Beyond these thermal
and kinetic effects, two compositional factors outside our model could
further modify the 660: 1) water in ringwoodite broadens the post-spinel
reaction and raises its pressure (Ghosh et al., 2013; Muir et al.,
2021), promoting additional deepening in cold slabs in contrast to
water's shallowing effect at the 410 (Smyth, 1987; Smyth \& Frost, 2002;
{van der Meijde et al.}, 2003), while 2) basalt accumulation modifies
reflectivity without systematically affecting depth ({Goes et al.},
2022; {Yu et al.}, 2023).

The directional contrast between kinetic effects at the 410 and 660 in
cold slabs is fundamental. At the 410, the positive Clapeyron slope
produces upward displacement under quasi-equilibrium conditions.
Decreasing reaction rates first broaden and deepen the discontinuity,
then ultimately re-sharpen it at a depressed position, reversing its
displacement sign relative to the equilibrium prediction (cf.~circle and
triangle datapoints in Figure~\ref{fig:660-structure}). At the 660, both
the equilibrium prediction and all three kinetic regimes produce
downward displacement (Figure~\ref{fig:seismic-profiles}). Kinetics
amplify rather than reverse the equilibrium signal, and the transition
from one regime to the next is monotonic. The 660 therefore lacks the
seismically diagnostic sign-reversal that marks the onset of
ultra-sluggish conditions at the 410.

Jointly observing both discontinuities in the same slab column exploits
these different behaviors. A sharp, deeply displaced 410 paired with a
broad, deeply displaced 660 points to ultra-sluggish olivine
\(\Leftrightarrow\) wadsleyite kinetics coexisting with intermediate to
ultra-sluggish ringwoodite decomposition. A broad, deepened 410 paired
with a broad, deepened 660 is consistent with intermediate kinetic
inhibition at both transitions. A modestly displaced, narrow 410 paired
with a modestly displaced, narrow 660 suggests quasi-equilibrium
conditions across the full transition zone
(Figure~\ref{fig:seismic-profiles}). These joint patterns exploit the
different saturation behaviors of the interface-controlled and
diffusion-controlled reaction mechanisms, though they only narrow
(rather than eliminate) the ambiguity with thermal, compositional, and
rheological contributions discussed below.

\begin{figure}
\centering
\includegraphics[width=0.75\linewidth,height=\textheight,keepaspectratio,alt={Synthetic seismic velocity profiles extracted along representative depth profiles in plume (left) and slab (right) simulations across the full range of tested kinetic prefactors Z\_\textbackslash mathrm\{ri\}. The black line shows the background adiabatic reference profile. Colored bands show the envelope of profiles from slow (dark red, Z\_\textbackslash mathrm\{ri\} = 1.6 \textbackslash times 10\^{}\{-3\} mol\^{}2 J\^{}\{-2\} s\^{}\{-1\}) to fast (yellow, Z\_\textbackslash mathrm\{ri\} = 1.6 \textbackslash times 10\^{}3 mol\^{}2 J\^{}\{-2\} s\^{}\{-1\}) kinetics. In plumes, the 410 and 520 are depressed, and the 660 is uplifted by elevated temperatures (+T). The narrow colored band at the 660 indicates that its position is insensitive to Z\_\textbackslash mathrm\{ri\}. In slabs, the 410 and 520 are uplifted by cold temperatures (-T), but counter-shifted by sluggish kinetic conditions (-Z\_\textbackslash mathrm\{ol\}, -Z\_\textbackslash mathrm\{wd\}). The 660 is depressed by both cold temperatures and kinetic inhibition (-T, -Z\_\textbackslash mathrm\{ri\}), with the wider colored band reflecting stronger sensitivity to Z\_\textbackslash mathrm\{ri\}. Unlike the 410, temperature (-T) and kinetic (Z\_\textbackslash mathrm\{ri\}) effects displace the 660 in the same direction in cold environments.}]{../figs/seisimc_profiles/seismic-composition.png}
\caption{Synthetic seismic velocity profiles extracted along
representative depth profiles in plume (left) and slab (right)
simulations across the full range of tested kinetic prefactors
\(Z_\mathrm{ri}\). The black line shows the background adiabatic
reference profile. Colored bands show the envelope of profiles from slow
(dark red, \(Z_\mathrm{ri}\) = 1.6 \(\times\) 10\(^{-3}\) mol\(^2\)
J\(^{-2}\) s\(^{-1}\)) to fast (yellow, \(Z_\mathrm{ri}\) = 1.6
\(\times\) 10\(^3\) mol\(^2\) J\(^{-2}\) s\(^{-1}\)) kinetics. In
plumes, the 410 and 520 are depressed, and the 660 is uplifted by
elevated temperatures (\(+T\)). The narrow colored band at the 660
indicates that its position is insensitive to \(Z_\mathrm{ri}\). In
slabs, the 410 and 520 are uplifted by cold temperatures (\(-T\)), but
counter-shifted by sluggish kinetic conditions (\(-Z_\mathrm{ol}\),
\(-Z_\mathrm{wd}\)). The 660 is depressed by both cold temperatures and
kinetic inhibition (\(-T\), \(-Z_\mathrm{ri}\)), with the wider colored
band reflecting stronger sensitivity to \(Z_\mathrm{ri}\). Unlike the
410, temperature (\(-T\)) and kinetic (\(Z_\mathrm{ri}\)) effects
displace the 660 in the same direction in cold
environments.}\label{fig:seismic-profiles}
\end{figure}

\subsection{Constraining Kinetics from Seismic
Observations}\label{sec:constraining-kinetics-from-seismic-observations}

Seismic observations in several subduction zones show the 660 displaced
over a broad area while 410 displacement is either more localized or
undetected, possibly because narrow structures escape detection.
Examples include subduction zones beneath South America (Ba et al.,
2025), East Asia (Liu et al., 2022; Sun et al., 2020; Tang et al.,
2023), and Europe (Cottaar \& Deuss, 2016; Kalmár et al., 2025). These
studies generally invoke water or the akimotoite \(\Leftrightarrow\)
bridgmanite transition to explain the pattern, but our results show
kinetic inhibition offers an alternative explanation for at least part
of the signature.

Beyond overall displacement, a handful of studies also constrain the
width of the transition, offering an independent test of our
kinetic-inhibition hypothesis. Prior work often attributes such
broadening to the joint decomposition of ringwoodite and garnet, but our
results show that kinetic effects alone can reproduce the observed
widths. Inversion of receiver functions over a broad area beneath the US
shows a velocity gradient over 25 km (Schmandt, 2012). More locally,
frequency-dependent receiver function observations constrain a 30--40 km
wide 660 beneath Chile--Argentina (Bonatto et al., 2020), and
triplications show a 50--70 km transition beneath Northeast China (Wang
\& Niu, 2010); both locations are associated with subducted slabs. The
Chile--Argentina width falls within the range predicted by kinetic
effects in the stagnation regime (Figure~\ref{fig:660-structure-comp};
Table S1), and, following the trends discussed in
Section~\ref{sec:structure-of-the-660}, the 50--70 km transition beneath
Northeast China is consistent with ultra-sluggish ringwoodite
decomposition within a strong, quickly descending slab.

Beyond depth and width, the 660 also displays a structural complexity
that generally exceeds the 410. The discontinuity can appear split
(Boyce \& Cottaar, 2021; Deuss et al., 2006) or difficult to resolve in
long-period PP precursor data (Deuss, 2009; Lessing et al., 2015), and
observed amplitude variations in precursors (Waszek et al., 2021; {Yu et
al.}, 2023) and high-frequency scattering (Day \& Deuss, 2013; Wu et
al., 2019) point to compositional variations around the boundary. Some
of this apparent complexity may itself be a kinetic signature. In
higher-frequency data, such as receiver functions, a broad transition
may resemble a split one. Determining how much kinetic effects
contribute to these observations (versus genuine compositional
variation) requires forward modeling of synthetic seismic data (e.g.,
Lessing et al., 2015).

Seismic observations alone, however, cannot fully separate kinetic
effects from thermal, compositional, and rheological contributions. The
boundaries between regimes shift systematically with \(B\), inflow
velocity, and the viscosity jump factor \(b\), so the same 660 structure
is consistent with many kinetic-rheological combinations. As previously
stated, joint analysis of the 410 and 660 in the same slab column
narrows, rather than eliminates, this ambiguity. However, two laboratory
priorities would translate this framework into quantitative
observational constraints. The first is high-pressure experiments that
reduce uncertainties in ringwoodite decomposition kinetic parameters
under cold slab conditions with small pressure oversteps. The second is
explicit kinetic models for the reverse post-spinel reaction (e.g.,
Lessing et al., 2022).

\subsection{Slab Dynamics and the Diversity of Subduction
Behavior}\label{sec:slab-dynamics-and-the-diversity-of-subduction-behavior}

Ringwoodite decomposition kinetics independently promote slab stagnation
through buoyancy accumulation in the metastable ringwoodite zone,
without requiring trench retreat or a prescribed lower-mantle viscosity
jump. Slabs progress from continuous penetration in the
quasi-equilibrium regime, through temporary stalling and flattening in
the intermediate regime, to full ponding above the 660 in the stagnation
regime (Figure~\ref{fig:slab-composition-660-b01};
Section~\ref{sec:structure-of-the-660}). These results complement the
equilibrium stagnation mechanisms identified by Garel et al. (2014),
{Agrusta et al.} (2017), and {Goes et al.} (2017). Where those studies
identify the Clapeyron slope, lower-mantle viscosity increase, and
trench retreat as necessary joint conditions for stagnation, our results
establish kinetic inhibition as an additional, independently variable
impedance. A lower-mantle viscosity jump reinforces kinetic stagnation
in our simulations (Figures S4--S7), confirming that kinetic and viscous
impedances at the 660 are additive.

The diversity of slab behavior documented globally by Fukao \& Obayashi
(2013) maps onto this three-regime framework. Stagnant western Pacific
slabs with 660 depressions of 30--60 km (Jiang et al., 2015; Tauzin et
al., 2017) are consistent with the intermediate regime at moderate
\(Z_\mathrm{ri}\), particularly where faster descent velocity through
the 660 amplifies metastable accumulation. The deepest observed
depressions exceeding 60 km require either stagnation regime kinetics,
the akimotoite \(\Leftrightarrow\) bridgmanite contribution ({Chanyshev
et al.}, 2022), or both acting together. Steeply penetrating slabs are
consistent with the quasi-equilibrium regime. Slabs that deflect at the
660 before eventually sinking may record temporal evolution within the
intermediate or stagnation regimes, where metastable ringwoodite is
gradually consumed until buoyancy is lost and the slab resumes descent,
consistent with Mao \& Zhong (2018) who found that stagnation is
transient on timescales of 20--30 Ma. A separate geodynamic modeling
study of ringwoodite decomposition kinetics predicts slab residence
times at the base of the transition zone of 150--160 Ma before buoyancy
is lost and penetration resumes (Chapman \& Clarke, 2021), a timescale
consistent with the persistent metastable zones in our stagnation regime
and with the long-lived sub-horizontal ponding documented in the western
Pacific (Fukao \& Obayashi, 2013).

Descent rates offer a further point of comparison with nature, though a
more demanding one given the inherent limitations of our
phenomenological model. Descent rates through the 660 in the
quasi-equilibrium regime fall somewhat below the natural interquartile
range of 2.6--4.8 cm yr\(^{-1}\) (Lallemand et al., 2005), a discrepancy
also present in the 410 results of Kerswell et al. (2026). Our model
already includes slab pull, the dominant driving force in nature, but
omits other controls on descent rate, including upper mantle viscosity
structure, chemical buoyancy contrasts between basaltic and harzburgitic
layers, volatile content, and 2D geometric simplifications that neglect
trench retreat and lateral flow. Any of these factors could shift
simulated rates closer to natural values. Because our simulations are
designed to isolate and demonstrate kinetic effects rather than to
reproduce the full boundary conditions of natural subduction, the
absolute descent rates are not expected to match observations
quantitatively. The qualitative correspondence with observed slab
diversity discussed above is the more direct test of our
kinetic-inhibition hypothesis. Incorporating the omitted physics would
be needed to translate our results into quantitative natural descent
rates, and represents a natural extension of this framework.

Despite all of these considerations, the clearest result of our
simulations is that kinetic inhibition does not act in isolation, but
jointly with rheology and kinematics to control subduction dynamics at
the 660. Stronger slabs (\(B\) = 10) descend faster through the 660 due
to their structural coherence and produce smaller displacements in the
quasi-equilibrium regime, where kinetics are fast enough to keep pace
with the slab. In the intermediate and stagnation regimes, however,
their faster descent through the 660 carries more material into the
metastable zone, amplifying 660 depression relative to weaker slabs.
Weaker slabs with faster descent rates through the 660 generally amplify
metastability and shift both regime boundaries to higher
\(Z_\mathrm{ri}\) (Figure~\ref{fig:660-structure-comp}). Small changes
in convergence rate or rheological strength near a regime boundary could
therefore drive episodic transitions between penetrating and stagnant
behavior at individual subduction zones over geological time ({Agrusta
et al.}, 2017; Kerswell et al., 2026).

\section{Conclusions}\label{sec:conclusions}

We couple diffusion-controlled ringwoodite decomposition kinetics to
compressible mantle flow simulations that also resolve the 410 and 520
transitions, varying the kinetic prefactor across six orders of
magnitude alongside slab strength, mantle viscosity structure, and
inflow velocity. Our simulations show that the different reaction
mechanisms governing the 410 and 660 produce fundamentally different
controls on discontinuity structure and slab dynamics across the MTZ.

Ringwoodite decomposition kinetics exert first-order control on slab
dynamics at the 660, modulated by mantle viscosity structure, descent
velocity, and slab strength. Slabs progress through three regimes as
reaction rates slow. Near equilibrium, slabs descend continuously
through the 660. As metastable ringwoodite accumulates, descent
progressively slows and the 660 broadens and deepens. At the slowest
reaction rates, slabs fully stagnate above the 660, descent rate drops
by an order of magnitude, and 660 depression reaches two to four times
the equilibrium prediction. This kinetic inhibition independently
promotes slab stagnation, complementing and amplifying the established
equilibrium mechanisms of the endothermic Clapeyron slope and
lower-mantle viscosity jump. The resulting seismic signatures differ and
become diagnostic. Quasi-equilibrium kinetics produce narrow, modestly
depressed reflectors, while stagnation produces broad, deeply depressed
reflectors with reduced amplitude, offering an alternative explanation
for observations often attributed instead to akimotoite, garnet, or
water.

The character of these regimes differs fundamentally from those at the
410. There, the interface-controlled driving force saturates at high
overstepping, permitting abrupt re-sharpening once temperatures overcome
the kinetic barrier. At the 660, the quadratic diffusion-controlled
driving force grows continuously with overstepping, so displacement and
width increase monotonically from quasi-equilibrium through stagnation
with no re-sharpening at any tested condition. Kinetic stagnation at the
660 is therefore progressive and persistent rather than
threshold-controlled, a distinction essential for correctly interpreting
the diverse seismic signatures of the MTZ.

Realizing the observational potential of this framework requires
progress on two fronts. High-pressure experiments targeting cold slab
conditions, including small pressure oversteps, would reduce current
uncertainties in activation energy and the kinetic prefactor, alongside
dedicated kinetic models for the reverse post-spinel reaction.
Frequency-dependent inversions that jointly constrain 660 topography and
width in well-studied subduction zones would help discriminate kinetic
from competing thermal and compositional explanations. Together with our
companion study of the 410, these advances would extend seismological
constraints to both the interface-controlled and diffusion-controlled
kinetic regimes of the MTZ.

\clearpage

\section*{Acknowledgements}\label{sec:acknowledgements}
\addcontentsline{toc}{section}{Acknowledgements}

This work was funded by the UKRI NERC Large Grant no. NE/V018477/1. All
computations were undertaken on Barkla2, part of the High Performance
Computing facilities at the University of Liverpool, who graciously
provided expert support. We thank the Computational Infrastructure for
Geodynamics (\url{https://geodynamics.org}) which is funded by the
National Science Foundation under awards EAR-0949446 and EAR-1550901 for
supporting the development of ASPECT.

\section*{Data Availability}\label{sec:data-availability}
\addcontentsline{toc}{section}{Data Availability}

All data, code, and relevant information for reproducing this work are
archived on the OSF (Kerswell, 2026a) and Zenodo (Kerswell, 2026b)
repositories. All code within these repositories is MIT Licensed and
free for use and distribution (see license details). ASPECT version
3.0.0 (Bangerth et al., 2024) was used for the computations in this
study and is freely available under the GPL v2.0 or later license.

\section*{Conflict of Interest}\label{sec:conflict-of-interest}
\addcontentsline{toc}{section}{Conflict of Interest}

The authors declare there are no conflicts of interest for this
manuscript.

\clearpage

\section*{References}\label{sec:references}
\addcontentsline{toc}{section}{References}

\protect\phantomsection\label{refs}
\begin{CSLReferences}{1}{0}
\bibitem[\citeproctext]{ref-agrusta2017}
{Agrusta, R., Goes, S., \& van Hunen, J.} (2017). Subducting-slab
transition-zone interaction: Stagnation, penetration and mode switches.
\emph{Earth and Planetary Science Letters}, \emph{464}, 10--23.

\bibitem[\citeproctext]{ref-ba2025}
Ba, K., Huang, J., Yu, Y., Song, J., Li, S., Shen, C., \& Liu, L.
(2025). Break-off and stagnation of the subducting {Nazca} slab in the
mantle transition zone beneath {West-Central South America}.
\emph{Tectonophysics}, \emph{911}, 230834.

\bibitem[\citeproctext]{ref-aspectmanual}
Bangerth, W., Dannberg, J., Fraters, M., Gassmöller, R., Glerum, A.,
Heister, T., et al. (2018, May). {ASPECT: Advanced Solver for Planetary
Evolution, Convection, and Tectonics {[}User Manual{]}} (Version 12).
Figshare. \url{https://doi.org/10.6084/m9.figshare.4865333}

\bibitem[\citeproctext]{ref-aspect-doi-v3.0.0}
Bangerth, W., Dannberg, J., Fraters, M., Gassmöller, R., Glerum, A.,
Heister, T., et al. (2024, December). {ASPECT: Advanced Solver for
Planetary Evolution, Convection, and Tectonics {[}Software{]}} (Version
3.0.0). Zenodo. \url{https://doi.org/10.5281/zenodo.14371679}

\bibitem[\citeproctext]{ref-bina1994}
Bina, C., \& Helffrich, G. (1994). Phase transition {Clapeyron} slopes
and transition zone seismic discontinuity topography. \emph{Journal of
Geophysical Research: Solid Earth}, \emph{99}(B8), 15853--15860.

\bibitem[\citeproctext]{ref-bonatto2020}
Bonatto, L., Piromallo, C., \& Cottaar, S. (2020). The transition zone
beneath {West Argentina}-{Central Chile} using {P-to-S} converted waves.
\emph{Journal of Geophysical Research: Solid Earth}, \emph{125}(8),
e2020JB019446.

\bibitem[\citeproctext]{ref-boyce2021}
Boyce, A., \& Cottaar, S. (2021). Insights into deep mantle
thermochemical contributions to {African} magmatism from converted
seismic phases. \emph{Geochemistry, Geophysics, Geosystems},
\emph{22}(3), e2020GC009478.

\bibitem[\citeproctext]{ref-burke1965}
Burke, J. (1965). The kinetics of phase transformations in metals.

\bibitem[\citeproctext]{ref-cahn1956}
Cahn, J. (1956). The kinetics of grain boundary nucleated reactions.
\emph{Acta Metallurgica}, \emph{4}(5), 449--459.

\bibitem[\citeproctext]{ref-chanyshev2022}
{Chanyshev, A., Ishii, T., Bondar, D., Bhat, S., Kim, E., Farla, R., et
al.} (2022). Depressed 660-km discontinuity caused by
akimotoite--bridgmanite transition. \emph{Nature}, \emph{601}(7891),
69--73.

\bibitem[\citeproctext]{ref-chapman2021}
Chapman, T., \& Clarke, G. (2021). Delayed growth of ferropericlase and
bridgmanite controls slab residence at the 660-km discontinuity.
\emph{Journal of Geophysical Research: Solid Earth}, \emph{126}(6),
e2020JB021487.

\bibitem[\citeproctext]{ref-clevenger2021}
Clevenger, T., \& Heister, T. (2021). Comparison between algebraic and
matrix-free geometric multigrid for a {Stokes} problem on adaptive
meshes with variable viscosity. \emph{Numerical Linear Algebra with
Applications}, \emph{28}(5), e2375.

\bibitem[\citeproctext]{ref-connolly2009}
Connolly, J. (2009). The geodynamic equation of state: What and how.
\emph{Geochemistry, Geophysics, Geosystems}, \emph{10}(10).

\bibitem[\citeproctext]{ref-cottaar2016}
Cottaar, S., \& Deuss, A. (2016). Large-scale mantle discontinuity
topography beneath {Europe}: Signature of akimotoite in subducting
slabs. \emph{Journal of Geophysical Research: Solid Earth},
\emph{121}(1), 279--292.

\bibitem[\citeproctext]{ref-cottaar2014}
Cottaar, S., Heister, T., Rose, I., \& Unterborn, C. (2014). {BurnMan}:
A lower mantle mineral physics toolkit. \emph{Geochemistry, Geophysics,
Geosystems}, \emph{15}(4), 1164--1179.

\bibitem[\citeproctext]{ref-day2013}
Day, E., \& Deuss, A. (2013). Reconciling {PP} and
{\(P^\prime P^\prime\)} precursor observations of a complex 660 km
seismic discontinuity. \emph{Geophysical Journal International},
\emph{194}(2), 834--838.

\bibitem[\citeproctext]{ref-deuss2009}
Deuss, A. (2009). Global observations of mantle discontinuities using
{SS} and {PP} precursors. \emph{Surveys in Geophysics}, \emph{30}(4),
301--326.

\bibitem[\citeproctext]{ref-deuss2006}
Deuss, A., Redfern, S., Chambers, K., \& Woodhouse, J. (2006). The
nature of the 660-kilometer discontinuity in {Earth}'s mantle from
global seismic observations of {PP} precursors. \emph{Science},
\emph{311}(5758), 198--201.

\bibitem[\citeproctext]{ref-dobson2014}
Dobson, D., \& Mariani, E. (2014). The kinetics of the reaction of
majorite plus ferropericlase to ringwoodite: Implications for mantle
upwellings crossing the 660 km discontinuity. \emph{Earth and Planetary
Science Letters}, \emph{408}, 110--118.

\bibitem[\citeproctext]{ref-flanagan1998}
Flanagan, M., \& Shearer, P. (1998). Global mapping of topography on
transition zone velocity discontinuities by stacking {SS} precursors.
\emph{Journal of Geophysical Research: Solid Earth}, \emph{103}(B2),
2673--2692.

\bibitem[\citeproctext]{ref-fraters2020}
Fraters, M. (2020, June). {The Geodynamic World Builder {[}Software{]}}
(Version 0.3.0). Zenodo. \url{https://doi.org/10.5281/zenodo.3900603}

\bibitem[\citeproctext]{ref-fraters2019}
{Fraters, M., Thieulot, C., van den Berg, A., \& Spakman, W.} (2019).
The {Geodynamic World Builder}: A solution for complex initial
conditions in numerical modeling. \emph{Solid Earth}, \emph{10}(5),
1785--1807.

\bibitem[\citeproctext]{ref-fukao2013}
Fukao, Y., \& Obayashi, M. (2013). Subducted slabs stagnant above,
penetrating through, and trapped below the 660 km discontinuity.
\emph{Journal of Geophysical Research: Solid Earth}, \emph{118}(11),
5920--5938.

\bibitem[\citeproctext]{ref-garel2014}
Garel, F., Goes, S., Davies, D., Davies, J., Kramer, S., \& Wilson, C.
(2014). Interaction of subducted slabs with the mantle transition-zone:
A regime diagram from 2-d thermo-mechanical models with a mobile trench
and an overriding plate. \emph{Geochemistry, Geophysics, Geosystems},
\emph{15}(5), 1739--1765.

\bibitem[\citeproctext]{ref-gassmoller2018}
Gassmöller, R., Lokavarapu, H., Heien, E., Puckett, E., \& Bangerth, W.
(2018). Flexible and scalable particle-in-cell methods with adaptive
mesh refinement for geodynamic computations. \emph{Geochemistry,
Geophysics, Geosystems}, \emph{19}(9), 3596--3604.

\bibitem[\citeproctext]{ref-gassmoller2020}
Gassmöller, R., Dannberg, J., Bangerth, W., Heister, T., \& Myhill, R.
(2020). On formulations of compressible mantle convection.
\emph{Geophysical Journal International}, \emph{221}(2), 1264--1280.

\bibitem[\citeproctext]{ref-ghosh2013}
Ghosh, S., Ohtani, E., Litasov, K., Suzuki, A., Dobson, D., \&
Funakoshi, K. (2013). Effect of water in depleted mantle on post-spinel
transition and implication for 660 km seismic discontinuity. \emph{Earth
and Planetary Science Letters}, \emph{371}, 103--111.

\bibitem[\citeproctext]{ref-goes2017}
{Goes, S., Agrusta, R., van Hunen, J., \& Garel, F.} (2017).
Subduction-transition zone interaction: A review. \emph{Geosphere},
\emph{13}(3), 644--664.

\bibitem[\citeproctext]{ref-goes2022}
{Goes, S., Yu, C., Ballmer, M., Yan, J., \& van der Hilst, R.} (2022).
Compositional heterogeneity in the mantle transition zone. \emph{Nature
Reviews Earth \& Environment}, \emph{3}(8), 533--550.

\bibitem[\citeproctext]{ref-green1979}
Green, D., Jaques, L., \& Hibberson, W. (1979). Petrogenesis of
mid-ocean ridge basalts. In \emph{The earth: Its origin, structure and
evolution} (pp. 265--300). Academic Press.

\bibitem[\citeproctext]{ref-heister2017}
Heister, T., Dannberg, J., Gassmöller, R., \& Bangerth, W. (2017). High
accuracy mantle convection simulation through modern numerical
methods--{II}: Realistic models and problems. \emph{Geophysical Journal
International}, \emph{210}(2), 833--851.

\bibitem[\citeproctext]{ref-hindmarsh2005}
Hindmarsh, A., Brown, P., Grant, K., Lee, S., Serban, R., Shumaker, D.,
\& Woodward, C. (2005). {SUNDIALS}: Suite of nonlinear and
differential/algebraic equation solvers. \emph{ACM Transactions on
Mathematical Software (TOMS)}, \emph{31}(3), 363--396.

\bibitem[\citeproctext]{ref-hirth2003}
Hirth, G., \& Kohlstedf, D. (2003). Rheology of the upper mantle and the
mantle wedge: A view from the experimentalists. \emph{Geophysical
Monograph-American Geophysical Union}, \emph{138}, 83--106.

\bibitem[\citeproctext]{ref-hosoya2005}
Hosoya, T., Kubo, T., Ohtani, E., Sano, A., \& Funakoshi, K. (2005).
Water controls the fields of metastable olivine in cold subducting
slabs. \emph{Geophysical Research Letters}, \emph{32}(17).

\bibitem[\citeproctext]{ref-ishii2019}
{Ishii, T., Huang, R., Myhill, R., Fei, H., Koemets, I., Liu, Z., et
al.} (2019). Sharp 660-km discontinuity controlled by extremely narrow
binary post-spinel transition. \emph{Nature Geoscience}, \emph{12}(10),
869--872.

\bibitem[\citeproctext]{ref-ita1992}
Ita, J., \& Stixrude, L. (1992). Petrology, elasticity, and composition
of the mantle transition zone. \emph{Journal of Geophysical Research:
Solid Earth}, \emph{97}(B5), 6849--6866.

\bibitem[\citeproctext]{ref-ji2004}
Ji, S. (2004). A generalized mixture rule for estimating the viscosity
of solid-liquid suspensions and mechanical properties of polyphase rocks
and composite materials. \emph{Journal of Geophysical Research: Solid
Earth}, \emph{109}(B10).

\bibitem[\citeproctext]{ref-jiang2015}
Jiang, G., Zhao, D., \& Zhang, G. (2015). Detection of metastable
olivine wedge in the western {Pacific} slab and its geodynamic
implications. \emph{Physics of the Earth and Planetary Interiors},
\emph{238}, 1--7.

\bibitem[\citeproctext]{ref-kalmar2025}
Kalmár, D., Petrescu, L., Hetényi, G., Michailos, K., Süle, B., Neagoe,
C., \& Bokelmann, G. (2025). Mantle transition zone analysis using
{P-to-S} receiver functions in the {Alpine--Carpathian--Dinarides}
region: Impact of plumes and slabs. \emph{Geophysical Journal
International}, \emph{243}(1), ggaf313.

\bibitem[\citeproctext]{ref-karato1984}
Karato, S. (1984). Grain-size distribution and rheology of the upper
mantle. \emph{Tectonophysics}, \emph{104}(1-2), 155--176.

\bibitem[\citeproctext]{ref-karato2008}
Karato, S. (2008). Deformation of earth materials. \emph{An Introduction
to the Rheology of Solid Earth}, \emph{463}.

\bibitem[\citeproctext]{ref-kerswell2026b}
Kerswell, B. (2026a, August).
{buchanankerswell/kerswell\_et\_al\_660\_kinetics: Publication Release
{[}Dataset{]}} (Version 1.0.0). OSF.
\url{https://doi.org/10.17605/OSF.IO/KUR93}

\bibitem[\citeproctext]{ref-kerswell2026c}
Kerswell, B. (2026b, August).
{buchanankerswell/kerswell\_et\_al\_660\_kinetics: Publication Release
{[}Software{]}} (Version 1.0.0). Zenodo.
\url{https://doi.org/xx.xxxx/zenodo.xxxxxxxx}

\bibitem[\citeproctext]{ref-kerswell2026a}
Kerswell, B., Wheeler, J., Gassmöller, R., Davies, J., Papanagnou, I.,
\& Cottaar, S. (2026). Beyond equilibrium: Kinetic thresholds and
rheological feedbacks create a potentially complex 410 in slab regions.
\emph{Journal of Geophysical Research: Solid Earth}, \emph{131}(7),
e2026JB033781.

\bibitem[\citeproctext]{ref-kronbichler2012}
Kronbichler, M., Heister, T., \& Bangerth, W. (2012). High accuracy
mantle convection simulation through modern numerical methods.
\emph{Geophysical Journal International}, \emph{191}(1), 12--29.

\bibitem[\citeproctext]{ref-kubo2000}
Kubo, T., Ohtani, E., Kato, T., Urakawa, S., Suzuki, A., Kanbe, Y., et
al. (2000). Formation of metastable assemblages and mechanisms of the
grain-size reduction in the postspinel transformation of
{Mg\(_2\)SiO\(_4\)}. \emph{Geophysical Research Letters}, \emph{27}(6),
807--810.

\bibitem[\citeproctext]{ref-kubo2002}
Kubo, T., Ohtani, E., Kato, T., Urakawa, S., Suzuki, A., Kanbe, Y., et
al. (2002). Mechanisms and kinetics of the post-spinel transformation in
{Mg\(_2\)SiO\(_4\)}. \emph{Physics of the Earth and Planetary
Interiors}, \emph{129}(1-2), 153--171.

\bibitem[\citeproctext]{ref-lallemand2005}
Lallemand, S., Heuret, A., \& Boutelier, D. (2005). On the relationships
between slab dip, back-arc stress, upper plate absolute motion, and
crustal nature in subduction zones. \emph{Geochemistry, Geophysics,
Geosystems}, \emph{6}(9).

\bibitem[\citeproctext]{ref-lebedev2002}
Lebedev, S., Chevrot, S., \& Hilst, R. van der. (2002). Seismic evidence
for olivine phase changes at the 410-and 660-kilometer discontinuities.
\emph{Science}, \emph{296}(5571), 1300--1302.

\bibitem[\citeproctext]{ref-lessing2015}
Lessing, S., Thomas, C., Saki, M., Schmerr, N., \& Vanacore, E. (2015).
On the difficulties of detecting {PP} precursors. \emph{Geophysical
Journal International}, \emph{201}(3), 1666--1681.

\bibitem[\citeproctext]{ref-lessing2022}
Lessing, S., Dobson, D., Rost, S., Cobden, L., \& Thomas, C. (2022).
Kinetic effects on the 660-km-phase transition in mantle upstreams and
seismological implications. \emph{Geophysical Journal International},
\emph{231}(2), 877--893.

\bibitem[\citeproctext]{ref-litasov2005}
Litasov, K., Ohtani, E., Sano, A., Suzuki, A., \& Funakoshi, K. (2005).
In situ x-ray diffraction study of post-spinel transformation in a
peridotite mantle: Implication for the 660-km discontinuity. \emph{Earth
and Planetary Science Letters}, \emph{238}(3-4), 311--328.

\bibitem[\citeproctext]{ref-liu2022}
Liu, L., Gao, S., Liu, K., Griffin, W., Li, S., Tong, S., \& Ning, J.
(2022). Mantle dynamics of the {North China Craton}: New insights from
mantle transition zone imaging constrained by {P-to-S} receiver
functions. \emph{Geophysical Journal International}, \emph{231}(1),
629--637.

\bibitem[\citeproctext]{ref-mao2018}
Mao, W., \& Zhong, S. (2018). Slab stagnation due to a reduced viscosity
layer beneath the mantle transition zone. \emph{Nature Geoscience},
\emph{11}(11), 876--881.

\bibitem[\citeproctext]{ref-mao2021}
Mao, W., \& Zhong, S. (2021). Formation of horizontally deflected slabs
in the mantle transition zone caused by spinel-to-post-spinel phase
transition, its associated grainsize reduction effects, and trench
retreat. \emph{Geophysical Research Letters}, \emph{48}(15),
e2021GL093679.

\bibitem[\citeproctext]{ref-ming2006}
Ming, L., Kim, Y., Uchida, T., Wang, Y., \& Rivers, M. (2006). In situ
x-ray diffraction study of phase transitions of {FeTiO\(_3\)} at high
pressures and temperatures using a large-volume press and synchrotron
radiation. \emph{American Mineralogist}, \emph{91}(1), 120--126.

\bibitem[\citeproctext]{ref-muir2021}
Muir, J., Zhang, F., \& Brodholt, J. (2021). The effect of water on the
post-spinel transition and evidence for extreme water contents at the
bottom of the transition zone. \emph{Earth and Planetary Science
Letters}, \emph{565}, 116909.

\bibitem[\citeproctext]{ref-myhill2023}
Myhill, R., Cottaar, S., Heister, T., Rose, I., Unterborn, C., Dannberg,
J., \& Gassmöller, R. (2023). {BurnMan}--a {Python} toolkit for
planetary geophysics, geochemistry and thermodynamics.

\bibitem[\citeproctext]{ref-navrotsky1998}
Navrotsky, A. (1998). Energetics and crystal chemical systematics among
ilmenite, lithium niobate, and perovskite structures. \emph{Chemistry of
Materials}, \emph{10}(10), 2787--2793.

\bibitem[\citeproctext]{ref-papanagnou2023}
Papanagnou, I., Schuberth, B., \& Thomas, C. (2023). Geodynamic
predictions of seismic structure and discontinuity topography of the
mantle transition zone. \emph{Geophysical Journal International},
\emph{234}(1), 355--378.

\bibitem[\citeproctext]{ref-reynolds2023}
Reynolds, D., Gardner, D., Woodward, C., \& Chinomona, R. (2023).
{ARKODE}: A flexible {IVP} solver infrastructure for one-step methods.
\emph{ACM Transactions on Mathematical Software}, \emph{49}(2), 1--26.

\bibitem[\citeproctext]{ref-ringwood1991}
Ringwood, A. (1991). Phase transformations and their bearing on the
constitution and dynamics of the mantle. \emph{Geochimica Et
Cosmochimica Acta}, \emph{55}(8), 2083--2110.

\bibitem[\citeproctext]{ref-schmandt2012}
Schmandt, B. (2012). Mantle transition zone shear velocity gradients
beneath {USArray}. \emph{Earth and Planetary Science Letters},
\emph{355}, 119--130.

\bibitem[\citeproctext]{ref-sime2024}
{Sime, N., Wilson, C., \& van Keken, P.} (2024). Thermal modeling of
subduction zones with prescribed and evolving 2D and 3D slab geometries.
\emph{Progress in Earth and Planetary Science}, \emph{11}(1), 14.

\bibitem[\citeproctext]{ref-smyth1987}
Smyth, J. (1987). {\(\beta\)-Mg\(_2\)SiO\(_4\)}; a potential host for
water in the mantle? \emph{American Mineralogist}, \emph{72}(11-12),
1051--1055.

\bibitem[\citeproctext]{ref-smyth2002}
Smyth, J., \& Frost, D. (2002). The effect of water on the 410-km
discontinuity: An experimental study. \emph{Geophysical Research
Letters}, \emph{29}(10), 123--1.

\bibitem[\citeproctext]{ref-stixrude2022}
Stixrude, L., \& Lithgow-Bertelloni, C. (2022). Thermal expansivity,
heat capacity and bulk modulus of the mantle. \emph{Geophysical Journal
International}, \emph{228}(2), 1119--1149.

\bibitem[\citeproctext]{ref-sun2020}
Sun, M., Gao, S., Liu, K., \& Fu, X. (2020). Upper mantle and mantle
transition zone thermal and water content anomalies beneath {NE Asia}:
Constraints from receiver function imaging of the 410 and 660 km
discontinuities. \emph{Earth and Planetary Science Letters}, \emph{532},
116040.

\bibitem[\citeproctext]{ref-tang2014}
Tang, Y., Obayashi, M., Niu, F., Grand, S., Chen, Y., Kawakatsu, H., et
al. (2014). {Changbaishan} volcanism in northeast {China} linked to
subduction-induced mantle upwelling. \emph{Nature Geoscience},
\emph{7}(6), 470--475.

\bibitem[\citeproctext]{ref-tang2023}
Tang, Z., Julià, J., Mooney, W., Mai, P., Yu, H., \& Wu, Y. (2023).
Seismic image of the mantle transition zone beneath northeastern
{China}: Evidence for stagnant {Pacific} subducting slab, lithospheric
delamination and mantle upwelling. \emph{Geophysical Journal
International}, \emph{235}(2), 1872--1887.

\bibitem[\citeproctext]{ref-tauzin2017}
Tauzin, B., Kim, S., \& Kennett, B. (2017). Pervasive seismic
low-velocity zones within stagnant plates in the mantle transition zone:
Thermal or compositional origin? \emph{Earth and Planetary Science
Letters}, \emph{477}, 1--13.

\bibitem[\citeproctext]{ref-vandermeijde2003}
{van der Meijde, M., Marone, F., Giardini, D., \& van der Lee, S.}
(2003). Seismic evidence for water deep in {Earth}'s upper mantle.
\emph{Science}, \emph{300}(5625), 1556--1558.

\bibitem[\citeproctext]{ref-wang2010}
Wang, B., \& Niu, F. (2010). A broad 660 km discontinuity beneath
northeast {China} revealed by dense regional seismic networks in
{China}. \emph{Journal of Geophysical Research: Solid Earth},
\emph{115}(B6).

\bibitem[\citeproctext]{ref-waszek2021}
Waszek, L., Tauzin, B., Schmerr, N., Ballmer, M., \& Afonso, J. (2021).
A poorly mixed mantle transition zone and its thermal state inferred
from seismic waves. \emph{Nature Geoscience}, \emph{14}(12), 949--955.

\bibitem[\citeproctext]{ref-wheeler2015}
Wheeler, J. (2015a). Dramatic effects of stress on metamorphic
reactions: reply. \emph{Geology}, \emph{43}(11), e373--e373.

\bibitem[\citeproctext]{ref-wheeler2015b}
Wheeler, J. (2015b). Dramatic effects of stress on metamorphic
reactions: reply. \emph{Geology}, \emph{43}(11), e373--e373.

\bibitem[\citeproctext]{ref-wheeler2020}
Wheeler, J. (2020). A unifying basis for the interplay of stress and
chemical processes in the {Earth}: Support from diverse experiments.
\emph{Contributions to Mineralogy and Petrology}, \emph{175}(12), 116.

\bibitem[\citeproctext]{ref-wu2019}
Wu, W., Ni, S., \& Irving, J. (2019). Inferring {Earth}'s discontinuous
chemical layering from the 660-kilometer boundary topography.
\emph{Science}, \emph{363}(6428), 736--740.

\bibitem[\citeproctext]{ref-xu2008}
Xu, W., Lithgow-Bertelloni, C., Stixrude, L., \& Ritsema, J. (2008). The
effect of bulk composition and temperature on mantle seismic structure.
\emph{Earth and Planetary Science Letters}, \emph{275}(1-2), 70--79.

\bibitem[\citeproctext]{ref-yamazaki2005}
Yamazaki, D., Inoue, T., Okamoto, M., \& Irifune, T. (2005). Grain
growth kinetics of ringwoodite and its implication for rheology of the
subducting slab. \emph{Earth and Planetary Science Letters},
\emph{236}(3-4), 871--881.

\bibitem[\citeproctext]{ref-yu2023}
{Yu, C., Goes, S., Day, E., \& van der Hilst, R.} (2023). Seismic
evidence for global basalt accumulation in the mantle transition zone.
\emph{Science Advances}, \emph{9}(22), eadg0095.

\bibitem[\citeproctext]{ref-zener1946}
Zener, C. (1946). Kinetics of the decomposition of austenite.
\emph{Trans. Aime}, \emph{167}, 550--595.

\end{CSLReferences}


\end{document}
